% !TEX root = sga4.5-en.tex
% ENGLISH EDITION -- TRANSLATED FROM FRENCH (AI-assisted, needs human review).

\chapter{Applications of the trace formula to exponential sums}\label{VI}
\chaptermark{Applications to exponential sums}










In this exposé, I explain how the trace formula makes it possible to compute
or to study various exponential sums and how, combined with the
Weil conjecture, it can make it possible to bound them.

The first two sections give a ``user's guide'' to these tools. Section 3 is an exposition,
in cohomological language, of Weil's results on sums in $1$ variable. Sections $4$ to $6$ form a
detailed study of Gauss and Jacobi sums -- including the
old and recent results of Weil on Hecke characters
defined by Jacobi sums. In section $7$, we study a
multi-variable generalisation of Kloosterman sums. Finally, in
section $8$, one will find some indications of other uses which have
been made or can be made of these methods. 










\section*{Notations}\label{VI:0}





\subsection{}\label{VI:0-1}

We use the notations of \hyperref[II]{Rapport}, paragraph \ref{II:1}. We
will often have to consider a finite extension $\dF_{q^n}\subset \dF$ of
$\dF_q$. We denote by a subscript $0$ an object over $\dF_q$, and by a subscript $1$
an object over $\dF_{q^n}$. Replacing a subscript $0$ by a subscript $1$ (resp.
deleting the subscript) means extending scalars to $\dF_{q^n}$ (resp.
to $\dF$).





\subsection{}\label{VI:0-2}

We denote by $\ell$ a prime number $\ne p$. We freely use the
language of $\dQ_\ell$-sheaves, as well as that of $E_\lambda$-sheaves,
for $E_\lambda$ a finite extension of $\dQ_\ell$ (cf. \hyperref[II]{Rapport},
paragraph \ref{II:2} and especially \ref{II:2-11}). 





\subsection{}\label{VI:0-3}

Let $H^\bullet$ be a graded vector space. If $T$ is an endomorphism of
$H^\bullet$, we set (cf. \hyperref[IV]{Cycle}, \ref{IV:eq:1-3-5}) 
\[
  \tr(T,H^\bullet) = \sum (-1)^i \tr(T,H^i) \text{.}
\]










\section{Principles}\label{VI:1}





\subsection{}\label{VI:1-1}

Let $X_0$ be a separated scheme of finite type over $\dF_q$, $E_\lambda$ a
finite extension of $\dQ_\ell$ and $\sF_0$ an $E_\lambda$-sheaf on $X-0$. The
trace formula says that 
\begin{equation*}\tag{1.1.1}\label{VI:eq:1-1-1}
  \sum_{x\in X^F} \tr(F_x^\ast,\sF) = \sum (-1)^i \tr\left(F^\ast,\h_c^i(X,\sF)\right) \text{.}
\end{equation*}

With the notation \ref{VI:0-3}, the right-hand side is simply written
$\tr(F^\ast,\h_c^\bullet(X,\sF))$.

This trace formula for $E_\lambda$-sheaves can either be
deduced from \hyperref[II]{Rapport} \ref{II:4-10} by passage to the limit (cf.
\hyperref[II]{Rapport} \ref{II:4-11} to \ref{II:4-13}), or be deduced
from the trace formula for $\dQ_\ell$-sheaves (\hyperref[II]{Rapport}
\ref{II:3-2}) by the method of (\hyperref[III]{Fonctions $L$ mod. $\ell^n$},
\ref{III:4-3}).

We are going to interpret various exponential sums as the left-hand side
of \eqref{VI:eq:1-1-1}, for suitable $X_0$ and $\sF_0$. 





% NOTE: in this section, following the source (but not convention in the 
% LaTeX version, torsors are denoted in roman (not script) letters
\subsection{}\label{VI:1-2}

Let $A$ be a finite group. An \emph{$A$-torsor} on a scheme $X$ is a
sheaf $T$ on $X$, equipped with a right action of $A$, which, locally
(for the étale topology) on $X$, is isomorphic to the constant sheaf $A$ on
which $A$ acts by right translation.

If $\tau:A\to B$ is a homomorphism, and $T$ an $A$-torsor, there exists up to
unique isomorphism one and only one $B$-torsor $\tau(T)$, equipped with a
morphism of sheaves $\tau:T\to \tau(T)$ such that 
\begin{equation*}\tag{1.2.1}\label{VI:eq:1-2-1}
  \tau(t a) = \tau(t) \tau(a)\text{.}
\end{equation*}

If $T$ is an $A$-torsor on $X$, and $\rho$ a linear representation
of $A$: $\rho:A\to \gl(V)$ ($V$ a finite-dimensional vector space over
$E_\lambda$), there exists up to unique isomorphism a single
$E_\lambda$-sheaf $\sF$, smooth of rank $\dim(V)$, equipped with a morphism of
sheaves 
\begin{equation*}\tag{1.2.2}\label{VI:eq:1-2-2}
  \rho:T\to \underline{\operatorname{Isom}}(V,\sF) \text{.}
\end{equation*}
such that $\rho(t a) = \rho(t) \rho(a)$. We denote it $\rho(T)$. For
$\tau:A\to B$ and $\rho$ a representation of $B$, we have a canonical
isomorphism $\rho(\tau(T)) = (\rho\tau)(T)$.

In this section (except for the appendix), we consider only the case where $A$
is commutative and where $V=E_\lambda$: $\rho$ is a character
$A\to E_\lambda^\times$, and $\rho(T)$ a smooth $E_\lambda$-sheaf of rank
one. The morphism \eqref{VI:eq:1-2-2} is interpreted as an everywhere
non-zero morphism 
\begin{equation*}\tag{1.2.3}\label{VI:eq:1-2-3}
  \rho:T\to \rho(T) 
\end{equation*}
such that $\rho(t a) = \rho(a)\rho(t)$. 





\subsection{}\label{VI:1-3}

Let $S$ be a scheme, $G$ a commutative group scheme over $S$, and $G'$
an extension of commutative group schemes over $S$
\[\xymatrix{
  0 \ar[r] 
    & A \ar[r] 
    & G' \ar[r]^-\pi 
    & G \ar[r] 
    & 0 \text{.}
}\]
The sheaf $T$ of local sections of $\pi$ is then an $A$-torsor on
$G$.

If $X$ is a scheme over $S$, and $f,g$ two $S$-morphisms from $X$ to $G$,
since $T$ is defined by an extension, we have a canonical
isomorphism 
\begin{equation*}\tag{1.3.1}\label{VI:eq:1-3-1}
  (f+g)^\ast T = f^\ast T + g^\ast T \text{.}
\end{equation*}
On the left, $f+g$ is the sum of $f$ and $g$, in the sense of the group law of $G$:
on the right, $+$ denotes a sum of torsors.

If $f:X\to G$ factors through $G'$, $f^\ast T$ is trivial (and a
factorisation gives a trivialisation); up to (non-unique) isomorphism,
the $A$-torsor $f^\ast T$ depends only on the image of $f$ in
$\hom_S(X,G)/\pi\hom_S(X,G')$: it is given by the morphism $\partial$ in
the exact sequence 
\[\xymatrix{
  \hom(X,G') \ar[r] 
    & \hom(X,G) \ar[r]^-\partial 
    & \h^1(X,A) \text{.}
}\]

We will apply these constructions to the Kummer exact sequences
$0 \to \dmu_n \to \dG_m \to \dG_m \to 0$, and to the Lang torsors. 





\subsection{The Lang torsor}\label{VI:1-4}

Let $G_0$ be a connected commutative algebraic group over $\dF_q$ (for the
non-commutative case, see the appendix). The group law is denoted
\emph{multiplicatively}. The Lang isogeny 
\[
\fL:G_0 \to G_0: x\mapsto F x\cdot x^{-1}
\]
is étale; its image, an open subgroup of $G_0$, can only be
$G_0$ itself; its kernel is the finite group $G_0(\dF_q)$. The \emph{Lang
torsor} $L$ is the $G_0(\dF_q)$-torsor on $G_0$ defined by the exact
sequence 
\begin{equation*}\tag{1.4.1}\label{VI:eq:1-4-1}
\xymatrix{
  0 \ar[r] 
    & G_0(\dF_q) \ar[r] 
    & G_0 \ar[r]^-\fL 
    & G_0 \ar[r] 
    & 0 \text{.}
}
\end{equation*}


\paragraph{Notation}
We still denote by $L$ and $\fL$ (or $L_{(q)}$ and $\fL_{(q)}$) the objects deduced
from $L$ and $\fL$ by extension of the base field. If necessary, we
will specify this by a subscript $0$ or $1$.


\paragraph{Examples}
For $G_0=\dG_a$ or $\dG_m$, the sequences \eqref{VI:eq:1-4-1} read 
\begin{align*}\tag{1.4.2}\label{VI:eq:1-4-2}
&\xymatrix{
  0 \ar[r] 
    & \dF_q \ar[r] 
    & \dG_a \ar[r]^-{x^q-x} 
    & \dG_a \ar[r] 
    & 0 
} \\
\tag{1.4.3}\label{VI:eq:1-4-3}
&\xymatrix{
  0 \ar[r] 
    & \dmu_{q-1} \ar[r] 
    & \dG_m \ar[r]^-{x^{q-1}} 
    & \dG_m \ar[r] 
    & 0
}
\end{align*}





\subsection{}\label{VI:1-5}

Let us compute the endomorphism $F^\ast$ of the fibre of the Lang torsor at
$\gamma\in G^F=G_0(\dF_q)$. If $g\in \fL^{-1}(\gamma)\subset G$, we have
$F g = F g\cdot g^{-1}\cdot g = \gamma g$. Hence (\hyperref[II]{Rapport},
\ref{II:1-2}) 
\begin{equation*}\tag{1.5.1}\label{VI:eq:1-5-1}
  \text{on $L_0(G_0)_\gamma\simeq \fL^{-1}(\gamma)$, $F^\ast$ is $g\mapsto g\gamma^{-1}$.}
\end{equation*}





\subsection{}\label{VI:1-6}

Over $\dF_{q^n}$, the identity $F_{(q^n)}=F_{(q)}^n$ between endomorphisms of
$G_1$ implies that $\fL_{(q^n)}=\fL_{(q)}\circ \prod_{i=0}^{n-1} F_{(q)}^i$.
On $G_0(\dF_{q^n})$, $F_{(q)}^i$ acts as the element
$x\mapsto x^{q^i}$ of $\gal(\dF_{q^n}/\dF_q)$. On $G_0(\dF_{q^n})$,
$\prod F_{(q)}^i$ is therefore the composite of the norm
$G_0(\dF_{q^n})\to G_0(\dF_q)$ and the inclusion
$G_0(\dF_q)\subset G_0(\dF_{q^n})$: the diagram 
\begin{equation*}\tag{1.6.1}\label{VI:eq:1-6-1}
\xymatrix{
  0 \ar[r] 
    & G_0(\dF_{q^n}) \ar[r] \ar[d]^-N 
    & G_1 \ar[r]^-{\fL_{(q^n)}} \ar[d]^-{\prod_{i=0}^{n-1} F_{(q)}^i} 
    & G_1 \ar[r] \ar@{=}[d] 
    & 0 \\
  0 \ar[r] 
    & G_0(\dF_q) \ar[r] 
    & G_q \ar[r]^-{\fL_{(q)}} 
    & G_1 \ar[r] 
    & 0
}
\end{equation*}
is commutative. In the language of torsors, this provides a canonical
isomorphism between $G_0(\dF_q)$-torsors on $G_1$ 
\begin{equation*}\tag{1.6.2}\label{VI:eq:1-6-2}
  N L_{(q^n)} = L_{(q)} \text{.}
\end{equation*}





\begin{definition_}\label{VI:1-7}
Let $f_0:X_0\to G_0$ be a morphism and
$\chi:G_0(\dF_q)\to E_\lambda^\times$ a character. We set
$\sF(\chi,f_0) = \chi^{-1}(f_0^\ast (L_0(G_0))) = f_0^\ast\chi^{-1}(L_0(G_0))$.
\end{definition_}

We have the following functoriality properties:
\begin{enumerate}[a)]
  \item $\sF(\chi,f_0)$ is bimultiplicative in $\chi$ and $f_0$: 
    \begin{align*}\tag{1.7.1}\label{VI:eq:1-7-1}
      \sF(\chi,f_0'\cdot f_0'') &= \sF(\chi,f_0')\otimes \sF(\chi,f_0'') \text{,} \\
      \tag{1.7.2}\label{VI:eq:1-7-2}
      \sF(\chi'\chi'',\sF_0) &= \sF(\chi',f_0)\otimes \sF(\chi'',f_0) \text{.}
    \end{align*}
    This follows from the definition of $\chi(T)$, for $T$ a torsor, together
    with \eqref{VI:eq:1-3-1} for \eqref{VI:eq:1-7-1}.
  \item For $g_0:Y_0\to X_0$ a morphism, we have 
    \begin{equation*}\tag{1.7.3}\label{VI:eq:1-7-3}
      \sF(\chi,f_0\circ g_0) = g_0^\ast \sF(\chi,f_0) \text{.}
    \end{equation*}
    In particular, $\sF(\chi,f_0) = f_0^\ast \sF(\chi)$, where $\sF(\chi)$
    denotes $\sF(\chi,\operatorname{id}_{G_0})$,
  \item For $u_0:G_0\to H_0$ a morphism, and
    $\chi:H_0(\dF_q) \to E_\lambda^\times$, we have 
    \begin{equation*}\tag{1.7.4}\label{VI:eq:1-7-4}
      \sF(\chi,u_0 f_0) = \sF(\chi u_0,f_0) \text{.}
    \end{equation*}
  \item For $G_0=\prod_{i\in I} G_0^i$, $\chi$ of coordinates $\chi_i$
    ($i\in I$) and $f_0$ of coordinates $f_0^i$, it follows from
    \eqref{VI:eq:1-7-1} to \eqref{VI:eq:1-7-4} that 
    \begin{equation*}\tag{1.7.5}\label{VI:eq:1-7-5}
      \sF(\chi,f_0) = \bigotimes_{i\in I} \sF(\chi_i,f_0^i) \text{.}
    \end{equation*}    
\end{enumerate}

Note the $\chi^{-1}$ in definition \ref{VI:1-7}. It ensures that, for
$x\in X^F$, we have on $\sF(\chi,f_0)_x$ 
\begin{equation*}\tag{1.7.6}\label{VI:eq:1-7-6}
  F_x^\ast = \chi f_0(x)
\end{equation*}
(use that the fibre at $x$ of the morphism \eqref{VI:eq:1-2-2}, for
$\rho=\chi^{-1}$, commutes with $F_x^\ast$, and \eqref{VI:eq:1-5-1}).

If $\sF(\chi,f_0)_1$ is the sheaf deduced from $\sF(\chi,f_0)$ by extension
of the base field to $\dF_{q^n}$, we deduce from \eqref{VI:eq:1-6-2} that 
\begin{equation*}\tag{1.7.7}\label{VI:eq:1-7-7}
  \sF(\chi,f_0)_1 = \sF(\chi\circ N,f_1) \text{.}
\end{equation*}





\subsection{Abuse of notations}\label{VI:1-8}

(i) If $\Xi$ is a notation for the composite map
$\chi f_0:X_0(\dF_q) \to E_\lambda^\times$, we will sometimes write $\sF(\Xi)$ in
place of $\sF(\chi,f_0)$. Thanks to \eqref{VI:eq:1-7-1} to \eqref{VI:eq:1-7-5},
there is hardly any risk of ambiguity. For example:
\begin{enumerate}[a)]
  \item we write $\sF(\chi)$ for $\sF(\chi,\operatorname{id}_{G_0})$
    (notation already used in \ref{VI:1-7}b));
  \item we write $\sF(\chi f_0)$ for $\sF(\chi,f_0)$;
  \item with the notations of \ref{VI:1-7}d), we write
    $\sF(\prod \chi_i f_0^i)$ for $\sF(\chi,f))$.
\end{enumerate}

With this notation, \eqref{VI:eq:1-7-4} expresses that the expression
$\sF(\chi u_0 f_0)$ is unambiguous, \eqref{VI:eq:1-7-1} \eqref{VI:eq:1-7-2}
\eqref{VI:eq:1-7-5} express multiplicativity of $\sF(\Xi)$ in $\Xi$,
and \eqref{VI:eq:1-7-6} is rewritten $F_x^\ast = \Xi(x)$ on $\sF(\Xi)_x$.

(ii) If $X$ is a scheme over an extension $k$ of $\dF_q$, and $f$ is a
morphism from $X$ to $G_0\otimes_{\dF_q} k$, we still denote by $\sF(\chi,f)$,
$\sF(\chi f)$, or $\sF(\chi)$ (for $f$ an inclusion) the inverse image by
$f$ of the $E_\lambda$-sheaf deduced from $\chi^{-1}(L_0(G))$ by extension of
scalars from $\dF_q$ to $k$. If $f_0$ is the composite
$X\to G_0\otimes_{\dF_q} k \to G_0$, it is still the $\sF(\chi,f_0)$ of
\ref{VI:1-7}. For $k$ a finite extension of $\dF_q$, we have
$\sF(\chi f) = \chi(\chi N_{k/\dF_q} f)$.

Applying \eqref{VI:eq:1-1-1} to \eqref{VI:eq:1-7-6} and \eqref{VI:eq:1-7-7},
we find: 





% NOTE: originally a 'scholium'
\begin{theorem_}\label{VI:1-9}
Let $S_0$ be a separated scheme of finite type over $\dF_q$, $G_0$ a connected
commutative algebraic group over $\dF_q$, $f_0:S_0\to G_0$ a morphism and
$\chi:G_0(\dF_q) \to E_\lambda^\times$. We have 
\begin{equation*}\tag{1.9.1}\label{VI:eq:1-9-1}
  \sum_{s\in S_0(\dF_q)} \chi f_0(s) = \tr(F^\ast,\h_c^\bullet(S,\sF(\chi,f_0))) 
\end{equation*}
and, for every integer $n\geqslant 1$
\begin{equation*}\tag{1.9.2}\label{VI:eq:1-9-2}
  \sum_{s\in S_0(\dF_{q^n})} \chi\circ N_{\dF_{q^n}/\dF_q} f_0(s) = \tr({F^\ast}^n, \h_c^\bullet\left(S,\sF(\chi,f_0))\right) \text{.}
\end{equation*}
\end{theorem_}





\addtocounter{subsubsection}{2}
\subsubsection{Remark}\label{VI:1-9-3}

Take for $G_0$ a product. With the notations of \ref{VI:1-7}d) and
\ref{VI:1-8}c), formula \eqref{VI:eq:1-9-2} becomes 
\[
  \sum_{s\in S_0(\dF_{q^n})} \prod_i \chi_i\circ N_{\dF_{q^n}/\dF_q} (f_0^i(s)) = \tr\Big({F^\ast}^n, \h_c^\bullet\Big(S,\sF\Big(\prod_i \chi_i f_0^i\Big)\Big)\Big) \text{.}
\]





\subsubsection{Remark}\label{VI:1-9-4}

Take $G_0=\{e\}$. Formula \eqref{VI:eq:1-9-1} becomes 
\[
  \left|S_0(\dF_q)\right| = \sum_{s\in S_0(\dF_q)} 1 = \tr(F^\ast,\h_c^\bullet(S,\dQ_\ell)) \text{.}
\]

It is rare that one can explicitly compute the right-hand side of
\eqref{VI:eq:1-9-1}. Here is an amusing example, with $G_0=\{e\}$. 





\subsection{Example}\label{VI:1-10}

A non-singular projective quadric of odd dimension over \texorpdfstring{$\dF_q$}{Fq} has the same number of rational points as the projective space over \texorpdfstring{$\dF_q$}{Fq} of the same dimension.

If, over $\dC$, $X$ is a non-singular quadric hypersurface in
projective space $\dP^{2 N}$, and $Y$ a hyperplane, we know that for
$i\leqslant 2\dim(X)=2\dim(Y)$, the inclusions 
$X\hookrightarrow \dP^{2 N}\hookleftarrow Y$ induce isomorphisms (in
ordinary cohomology) 
\[\xymatrix{
  \h^i(X,\dQ) 
    & \ar[l]_-\sim \h^i(\dP^{2 N},\dQ) \ar[r]^-\sim 
    & \h^i(Y,\dQ) \text{.}
}\]

By specialisation, it follows that if $X_0'$ is a non-singular
quadric hypersurface in projective space $\dP_0^{2 N}$ over $\dF_q$,
and $Y_0'$ a hyperplane, the inclusions
$X_0'\hookrightarrow \dP_0^{2 N}\hookleftarrow Y_0'$ induce
isomorphisms 
\[\xymatrix{
  \h^i(X',\dQ_\ell) 
    & \ar[l]_-\sim \h^i(\dP^{2 N},\dQ_\ell) \ar[r]^-\sim 
    & \h^i(Y',\dQ_\ell) \text{.}
}\]
These isomorphisms commute with $F^\ast$, and we apply the trace formula. 





\subsection{}\label{VI:1-11}

One can also use the trace formula to understand the classical
formulas giving the number of rational points of linear groups over
finite fields, and those of certain homogeneous spaces (see section
8). 





\subsection{}\label{VI:1-12}

We have a dictionary making it possible to translate into cohomological terms
various types of classical manipulations on exponential
sums. This dictionary will be given in section 2. The
cohomological statement being more ``geometric,'' it can sometimes apply
to situations where the classical argument applies only after an
extension of the base finite field. Here is an example (a special case of a
theorem of Kazhdan). 





\begin{theorem_}[Kazhdan]\label{VI:1-13}
Let $X_0$ be a scheme over $\dF_q$. Suppose there exist an action $\rho$ of
$\dG_a$ on $X$, and a morphism $f:X\to Y$ of schemes over $\dF$, making
$X$ a $\dG_a$-torsor over $Y$. Then the number of rational points of $X_0$
is divisible by $q$.
\end{theorem_}

Suppose first that $\rho$, $Y$ and $f$ are defined over $\dF_q$.


\paragraph{Classical argument}
The fibres of $f:X_0(\dF_q)\to Y_0(\dF_q)$ all have $q$ elements:
$|X_0(\dF_q)| = q\cdot |Y_0(\dF_q)|$, whence the divisibility.


\paragraph{Cohomological translation}
The higher direct image sheaves $\eR^i f_!\dQ_\ell$ are 
\[
  \eR^i f_!\dQ_\ell = \begin{cases}
                       \dQ_\ell(-1) & \text{if $i = 2$} \\
                       0            & \text{for $i\ne 2$}
                     \end{cases}
\]
The Leray spectral sequence of $f$ therefore degenerates into an isomorphism 
\[
  \h_c^i(X,\dQ_\ell) = \h_c^{i-2}(Y,\dQ_\ell)(-1) \text{.}
\]
To understand the effect of a Tate twist on the eigenvalues of
Frobenius, the most convenient is to use the Galois point of view
(\hyperref[II]{Rapport} \ref{II:1-8}), and to note that the Frobenius
substitution $\varphi$ acts on $\dQ_\ell(1)$ by multiplication by $q$. The
eigenvalues of $F^\ast$ on $\h_c^p(X,\dQ_\ell)$ are therefore the products by
$q$ of the eigenvalues of $F^\ast$ acting on $\h_c^{i-2}(Y,\dQ_\ell)$. We
know that these are algebraic integers \cite[XXI 5.2.2]{sga7}; hence 

\begin{equation*}\tag{$*$}\label{VI:eq:*}
  \substack{\displaystyle\text{the eigenvalues of $F^\ast$ acting on $\h_c^i(X,\dQ_\ell)$ are} \\
  \displaystyle\text{algebraic integers divisible by $q$.}}
\end{equation*}


\paragraph{Descent}
Let us return to the hypotheses of the theorem. For suitable $n$, we may
suppose that $\rho$, $Y$ and $f$ are defined over $\dF_{q^n}$. The
Frobenius morphism relative to $\dF_{q^n}$ being the $n$-th power of that
relative to $\dF_q$, the statement \eqref{VI:eq:*} for the scheme over
$\dF_{q^n}$ deduced from $X_0$ by extension of scalars gives us: 
\begin{equation*}\tag{$**$}\label{VI:eq:**}
  \substack{\displaystyle\text{the $n$-th powers of the eigenvalues of $F^\ast$ acting on} \\ \displaystyle\text{$\h_c^i(X,\dQ_\ell)$ are algebraic integers divisible by $q^n$.}}
\end{equation*}

This assertion implies \eqref{VI:eq:*}. The number of rational points of
$X_0$, namely $\sum (-1)^i \tr(F^\ast,\h^i(X,\dQ_\ell))$, is therefore an
algebraic integer divisible by $q$. This implies that it is divisible by $q$ as
a rational integer. 





\subsection{}\label{VI:1-14}

Formula \ref{VI:1-9-3} controls the dependence on $n$ of the
exponential sum on the left-hand side. It implies identities between an
exponential sum and those deduced from it by ``extension of
scalars.'' The most famous of these identities is that of
Hasse-Davenport: let $\chi:\dF_q^\times \to \dC^\times$ and
$\psi:\dF_q\to \dC^\times$ be characters, $\psi$ non-trivial, and
let us define the Gauss sum $\tau(\chi,\psi)$ by 
\begin{equation*}\tag{1.14.1}\label{VI:eq:1-14-1}
  \tau(\chi,\psi) = - \sum_{x\in \dF_q^\times} \psi(x) \chi^{-1}(x) \text{.}
\end{equation*}





\begin{theorem_}[Hasse-Davenport]\label{VI:1-15}
We have 
\begin{equation*}\tag{1.15.1}\label{VI:eq:1-15-1}
  \tau\left(\chi\circ N_{\dF_{q^n}/\dF_q},\psi\circ \tr_{\dF_{q^n}/\dF_q}\right) = \tau(\chi,\psi)^n \text{.}
\end{equation*}
\end{theorem_}

Let $E\subset \dC$ be a number field containing the values of $\chi$ and
$\psi$, and $E_\lambda$ the completion of $E$ at a place of
characteristic $\ell\ne p$. Identity \eqref{VI:eq:1-15-1} is an identity in $E$. It
amounts to the same to prove it in $\dC$, or in $E_\lambda$. We will
prove it in $E_\lambda$, regarding $\chi$ and $\psi$ as taking values
in $E_\lambda^\times$.

We apply \ref{VI:1-9-3} with $X_0=\dG_m$ over $\dF_q$ and
$G_0=\dG_m\times \dG_a$. The $\h_c^i(\dG_m,\sF(\chi^{-1}\psi))$ vanish for
$i\ne 1$, and the $\h_c^1$ is of dimension $1$ (\ref{VI:4-2}). By
\ref{VI:1-9-3},
$\tau(\chi\circ N_{\dF_{q^n}/\dF_q},\psi\circ \tr_{\dF_{q^n}/\dF_q})$ is
the unique eigenvalue of $(F^\ast)^n$ on $\h_c^1$, and \eqref{VI:eq:1-15-1}
follows.

Analogous examples are given in Weil \cite[App V]{we74}. 





\subsection{}\label{VI:1-16}

An $E_\lambda$-sheaf $\sF_0$ on $X-0$ is said to be \emph{pointwise of
weight $n$} if for every closed point $x\in |X_0|$, the eigenvalues of
$F_x^\ast$ (\hyperref[II]{Rapport} \ref{II:1-2}) are algebraic numbers
all of whose complex conjugates have absolute value
$q_x^{n/2}$, where $q_x$ is the number of elements of $k(x)$. The following
theorem will be proved in \cite{de80}. 





\begin{theorem_}\label{VI:1-17}
If $\sF_0$ is pointwise of weight $n$, for every eigenvalue $\alpha$
of the endomorphism $F^\ast$ of $\h_c^i(X,\sF)$, there exists an integer
$m\leqslant n+i$ such that the complex conjugates of $\alpha$ are all of
absolute value $q^{m/2}$.
\end{theorem_}

The sheaves considered in \ref{VI:1-9} are of weight $0$. The theorem
therefore provides for the sum \ref{VI:eq:1-9-1} (or rather, for all its
complex conjugates) the bound 
\begin{equation*}\tag{1.17.1}\label{VI:eq:1-17-1}
  \left|\sum_{s\in S_0(\dF_q)} \prod_i \chi_i(f^i(s)) \right| \leqslant \sum_i \dim \h_c^i\left(S,\sF\left(\prod \chi_i f^i\right)\right)\cdot q^{i/2} \text{.}
\end{equation*}





\subsection{Remarks}\label{VI:1-18}

\begin{enumerate}[a)]
  \item If $\dim S=n$, for every $E_\lambda$-sheaf $\sF$ on $S$, we have 
    \[
      \h_c^i(S,\sF) = 0 \qquad \text{for $i\notin [0,2n]$.}
    \]
  \item The group $\h_c^0(X,\sF)$ is the group of global sections of $\sF$
    on $S$ with proper support. If $S$ is connected, non-complete and
    $\sF$ smooth, indeed if $S$ is connected, $\sF$ smooth of rank one, and
    non-constant, this group is zero.
  \item If $S$ is smooth purely of dimension $n$, and $\sF$ smooth,
    Poincaré duality says that the vector spaces
    $\h_c^i(X,\sF)$ and $\h^{2n-i}(X,\sF^\vee)(2n)$ are dual to each other
    (for the effect of a Tate twist on the eigenvalues of Frobenius,
    cf. the $2$nd part of the proof of \ref{VI:1-13}).
  \item If $\dim S=n$, Poincaré duality makes it possible to compute
    $\h_c^{2n}(S,\sF)$ as follows: we take an open $U$ of $S_\text{red}$, whose
    complement is of dimension $<n$, smooth purely of dimension $n$, and
    on which $\sF$ is smooth. We then have
    $\h_c^{2n}(U,\sF) \iso \h_c^{2n}(S,\sF)$, for in the long exact
    cohomology sequence, $\h_c^i(S\setminus U,\sF) = 0$ for $i=2n,2n-1$. By
    Poincaré, we therefore have 
    \[
      \h_c^{2n}(S,\sF) \iso \left(\h^0(U,\sF^\vee)(2n)\right)^\vee \text{.}
    \]
    
    Suppose $U$ connected, and let $u$ be a geometric point of $U$. The
    ``local system'' $\sF$ corresponds to a representation of
    $\pi_1(U,u)$ on $\sF_u$, and $\h^0(U,\sF^\vee)$ is
    $(\sF_u^\vee)^{\pi_1(U,u)}=((\sF_u)_{\pi_1(U,u)})^\vee$ (dual of the
    coinvariants). We therefore have 
    \[
      \h_c^{2n}(S,\sF) \simeq (\sF_u)_{\pi_1(U,u)}(-2n) \text{.}
    \]
  \item These remarks often allow the computation of $b_0$ and $b_{2n}$.
    Computing the other $b_i$ can be difficult. If $b_{2n}=0$ and
    $b_{2n-1}=0$, the bound \eqref{VI:eq:1-17-1} is a Lang-Weil-type bound,
    with, for $q$ large, a gain of $\sqrt q$ over the
    trivial bound in $O(q^n)$. One should however beware that the
    bound \eqref{VI:eq:1-17-1} does not formally imply the trivial
    bound 
    \[
      \left|\sum_{s\in S_0(\dF_q)} \prod_i \chi_i(f^i(s))\right| \leqslant \left| S_0(\dF_q)\right| \text{,}
    \]
    so it can be useful to take it into account (cf. the proof of \ref{VI:3-8}).
\end{enumerate}

I explain below a method to prove that $b_i=0$ for $i\ne n$.
When it applies, it provides an estimate in $O(q^{n/2})$. The
implicit constant in $O$ is $b_n=(-1)^n \chi(S,\sF(\prod \chi_i f^i))$ and
can be difficult to compute. 





\begin{proposition_}\label{VI:1-19}
Let $\bar X$ be a proper scheme over an algebraically closed field $k$,
$j:X\hookrightarrow bar X$ an open of $X$ and $\sF$ an $E_\lambda$-sheaf
on $X$. We suppose that
\begin{enumerate}[\indent a)]
  \item $(j_\ast\sF)_x = 0$ for $x\in \bar X\setminus X$, i.e.
    $j_!\sF\iso j_\ast\sF$;
  \item $\eR^i j_\ast \sF = 0$ for $i>0$.
\end{enumerate}
Then the maps 
\[
  \h_c^i(X,\sF) \to \h^i(X,\sF)
\]
are isomorphisms.
\end{proposition_}

Hypotheses a) b) mean that $j_! \sF \iso \eR j_\ast \sF$, whence 
\[
  \h_c^i(X,\sF) = \h^i(\bar X,j_!\sF) \iso \hh^i(\bar X,\eR j_\ast \sF) = \h^i(X,\sF) \text{.}
\]
In other words, the Leray spectral sequence for $j$ 
\[
  E_2^{pq} = \h^p(\bar X,\eR^q j_\ast \sF) \Rightarrow \h^{p+q}(X,\sF) 
\]
degenerates, by b) into isomorphisms 
\[
  \h^i(\bar X,j_\ast\sF) \iso \h^i(X,\sF) \text{,}
\]
while, by a) 
\[
  \h_c^i(X,\sF) = \h^i(\bar X,j_!\sF) \iso \h^i(\bar X,j_\ast\sF) \text{.}
\]





\subsubsection{Example}\label{VI:1-19-1}

If $X$ is a curve, or if $\bar X$ is smooth, $X$ the complement of a
normal-crossings divisor and $\sF$ tamely ramified,
hypothesis a) of \ref{VI:1-19} (no invariants under local monodromy)
implies hypothesis b). 





\begin{proposition_}\label{VI:1-20}
Let $X_0$ be a smooth affine scheme purely of dimension $n$ over $\dF_q$ and
$\sF_0$ a smooth $E_\lambda$-sheaf pointwise of weight $m$ on $X_0$. We
suppose that the maps $\h_c^i(X,\sF) \to \h^i(X,\sF)$ are
isomorphisms (cf. \ref{VI:1-17}). Then $\h_c^i(X,\sF) = 0$ for $i\ne n$, and
the complex conjugates of the eigenvalues of $F^\ast$ on $\h_c^n(X,\sF)$
are of absolute value $q^{(m+n)/2}$.
\end{proposition_}

Poincaré duality puts in duality $\h_c^i(X,\sF)$ (resp.
$\h_c^i(X,\sF^\vee)$) with $\h^{2n-i}(X,\sF^\vee(n))$ (resp.
$\h^{2n-i}(X,\sF(n))$). By \cite[XIV 3.2]{sga4}, the $\h^i$ without support
vanish for $i>n$. By duality, the $\h_c^i$ vanish for $i<n$; since
$\h_c^i(X,\sF) \iso \h^i(X,\sF)$, these groups vanish for $i\ne n$.
Apply \ref{VI:1-15} to $\sF_0$ and to $\sF_0^\vee(n)$ (of weight
$-m-2n$). We find that the complex conjugates $\alpha$ of the eigenvalues
of $F^\ast$ on $\h_c^n(X,\sF) = \h_c^n(X,\sF^\vee(n))^\vee$ satisfy 
\begin{align*}
  |\alpha| &\leqslant q^{(m+n)/2} && \text{and} \\
  |\alpha^{-1}| &\leqslant q^{((-m-2n)+n)/2} = q^{-(m+n)/2} \text{,}
\end{align*}
whence the assertion.




\subsection{Remark}\label{VI:1-21}

The last argument shows that if $X_0$ is a smooth separated scheme over
$\dF_q$, $\sF_0$ is a smooth $E_\lambda$-sheaf pointwise of weight
$m$ on $X_0$ and $\h_c^i(X,\sF) \hookrightarrow \h^i(X,\sF)$, then the
complex conjugates of the eigenvalues of $F^\ast$ on $\h_c^i(X,\sF)$ are
of absolute value $q^{(m+i)/2}$. 





\subsection*{Appendix. The Lang isogeny in the non-commutative case}




\subsection{}\label{VI:1-22}

Let $G-0$ be a connected algebraic group over $\dF_q$ and $\fL$ the morphism
$x\mapsto F x\cdot x^{-1}$ from $G-0$ to itself. It is a morphism of
$G_0$-homogeneous spaces, if at the source we let $G-0$ act by left
translation $x\mapsto (y\mapsto x y)$, at the target by
$x\mapsto (y\mapsto F x\cdot y\cdot x^{-1}$). We have $\fL(xy) = \fL(x)$ for
$y\in G_0(\dF_q)$, and $\fL$ induces an isomorphism $G_0/G_0(\dF_q)\iso G_0$.

As in the commutative case, $\fL$ therefore makes $G_0$ a $G_0(\dF_q)$-torsor
over $G_0$, the \emph{Lang torsor} $L$. If
$\rho:G_0(\dF_q) \to \gl(n,E_\lambda)$ is a linear representation of
$G_0(\dF_q)$, we will denote by $\sF(\rho)$ the $E_\lambda$-sheaf $\rho(L)$
(\ref{VI:1-2}). 





\subsection{}\label{VI:1-23}

Let $\gamma\in G^F=G_0(\dF_q)$, and let us compute the endomorphism $F^\ast$ of
$\sF(\rho)_\gamma$. Let $g\in G$ such that $\fL(g)=\gamma$, whence a frame
$\rho(g)\in \operatorname{Isom}(E_\lambda^n,\sF(\rho)_\gamma)$. For
$e\in E_\lambda^n$, we have 
\[
  F(\rho(g)(e)) = \rho(F g)(e) = \rho(\gamma g)(e) = \rho(g\cdot g^{-1}\gamma g)(e) \text{.}
\]
We will see that $g^{-1} \gamma g\in G_0(\dF_q)$, whence 
\[
  F(\rho(g)(e)) = \rho(g) \rho(g^{-1}\gamma g)(e)
\]
et 
\begin{equation*}\tag{1.23.1}\label{VI:eq:1-23-1}
  \rho(g)^{-1} (F_\gamma^\ast)^{-1} \rho(g) = \rho(g^{-1} \gamma g) \text{:}
\end{equation*}
$\rho(g)$ identifies the inverse of $F^\ast$ at $\gamma$ with the automorphism
$\rho(g^{-1} \gamma g)$ of $E_\lambda^n$. It remains to understand what
$g^{-1} \gamma g$ is. 





\begin{lemma_}\label{VI:1-24}
\begin{enumerate}[(i)]
  \item If $F g\cdot g^{-1} \in G_0(\dF_q)$, we have
    $g^{-1} (F g\cdot g^{-1}) g = g^{-1} F g\in G_0(\dF_q)$.
  \item Let $\gamma\in G_0(\dF_q)$, and $g$ such that
    $F g\cdot g^{-1} = \gamma$. The conjugacy class of $\gamma'=g^{-1} Fg$
    in $G_0(\dF_q)$ depends only on that of $\gamma$.
  \item The map induced by $\gamma\mapsto \gamma'$, from the set
    $G_0(\dF_q)^\natural$ of conjugacy classes of $G_0(\dF_q)$ to
    itself, is bijective. Its inverse is the map for the opposite
    group.
\end{enumerate}
\end{lemma_}

(i) If $\gamma= F g\cdot g^{-1}\in G_0(\dF_q)$, the identities $F g=\gamma g$
and
$F\gamma=\gamma$ give $F(g^{-1} F g) = (F g)^{-1} F F g = g^{-1} \gamma^{-1} F(\gamma g) = g^{-1} \gamma^{-1} F \gamma F g = g^{-1} F g$,
so that $g^{-1} F g \in G_0(\dF_q)$.

(ii) For $\alpha,\gamma\in G_0(\dF_q)$, the $g$ such that
$F g\cdot g^{-1} = \gamma$ form a right class under $G_0(\dF_q)$, and
if $F g\cdot g^{-1} = \gamma$, we have
$F(\alpha g \alpha^{-1})(\alpha g \alpha^{-1})^{-1} = \alpha \gamma\alpha^{-1}$.
For $\gamma$ running through a conjugacy class in $G_0(\dF_q)$, the $g$
such that $F g\cdot g^{-1} = \gamma$ therefore run through a double class under
$G_0(\dF_q)$, and $\gamma'= g^{-1} F g$ runs through a conjugacy class.

(iii) Finally, the inverse of $F g\cdot g^{-1}\mapsto g^{-1} F g$ is
$g^{-1} F g\mapsto F g\cdot g^{-1}$. 





\subsection{}\label{VI:1-25}

Formula \eqref{VI:eq:1-23-1} can be read, briefly: $F^\ast$ at
$\gamma$ is $\rho(\gamma')^{-1}$. If $\gamma$ belongs to the neutral
component of its centraliser, one may take $g$ in $Z(\gamma)$. We then have
$\gamma=\gamma'$. If this condition is not fulfilled, $\gamma$ and $\gamma'$
may fail to be conjugate. 





\subsection{An example}\label{VI:1-26}

For $G_0=\operatorname{SL}(2)$, $q$ odd,
$\alpha\in F_q^\times\setminus {\dF_q^\times}^2$ and $\bar\gamma$ the
conjugacy class of $-1(1+N)$, with $N^2=0$, $\bar\gamma'$ is the
conjugacy class of $-1\cdot (1+\alpha N)$, so that $\bar\gamma'\ne \bar\gamma$
if $N\ne 0$. 










\section{Dictionary}\label{VI:2}

An elementary manipulation on exponential sums can often
be seen as the reflection, via the trace formula, of a cohomological
statement. In this section, I list such statements, and their
reflection; they bear the same number, increased by a $\ast$ for
the cohomological statement. 





\subsection{}\label{VI:2-1}

``A sum of algebraic integers is an algebraic integer'' admits as
cohomological analogue the following theorem, proved in
\cite[XXI 5.2.2]{sga7} (cf. the use of \ref{VI:2-1}$^\ast$ in
\ref{VI:1-13}). 





\begin{theorem*}[2.1*]\label{VI:2-1*}
Let $X_0$ be a separated scheme of finite type over $\dF_q$ and $\sF_0$ an
$E_\lambda$-sheaf on $X_0$. If, for every $x\in |X_0|$, the eigenvalues
of $F_x^\ast$ on $\sF_0$ are algebraic integers, then the
eigenvalues of $F^\ast$ on $\h_c^i(X,\sF)$ are algebraic integers.
\end{theorem*}





\subsection{}\label{VI:2-2}

Other integrality results are proved in
\cite[XXI, 5.2.2 et 5.4]{sga7}: if $\dim(x_0)\leqslant n$, and $\alpha$ is
an eigenvalue of $F^\ast$ on $\h_c^i(X,\sF)$, we have:
\begin{enumerate}[a)]
  \item under the hypotheses of \ref{VI:2-1}$^\ast$, and if $i>n$, then
    $q^{i-n}$ divides $\alpha$;
  \item if the inverses of the eigenvalues of $F_x^\ast$ are algebraic
    integers, then $\alpha^{-1}$ is integral, except at $p$. More
    precisely, $q^{\inf(n,i)}\alpha^{-1}$ is an algebraic integer.
\end{enumerate}

For the sheaves considered in \ref{VI:1-9}, the eigenvalues of the
$F_x^\ast$ are roots of unity, so that the hypotheses of the
theorem and that of b) above are satisfied. 





\subsection{}\label{VI:2-3}

Let $f:A\to B$ be a map from one finite set to another, and
$\varepsilon$ a function on $A$. We have 
\begin{equation*}\tag{2.3.1}\label{VI:eq:2-3-1}
  \sum_{a\in A} \varepsilon(a) = \sum_{b\in B} \sum_{f(a)=b} \varepsilon(a) \text{.}
\end{equation*}





\subsection*{2.3*}\label{VI:2-3*}

Let $f:X\to Y$ be a morphism of separated schemes of finite type over $k$
algebraically closed, and $\sF$ an $E_\lambda$-sheaf on $X$. The
Leray spectral sequence of $f$ in proper-support cohomology reads 
\begin{align*}\tag{2.3.1*}\label{VI:eq:2-3-1*}
  E_2^{p q} = \h_c^p(Y,\eR^q f_! \sF) \Rightarrow \h_c^{p+q}(X,\sF) \text{.}
\end{align*}

Let $f_0:X_0 \to Y_0$ be a morphism of separated schemes of finite type over
$\dF_q$ and $\sF_0$ an $E_\lambda$-sheaf on $X_0$. Between
\eqref{VI:eq:2-3-1} and \eqref{VI:eq:2-3-1*} we have the following compatibility. 

\begin{enumerate}[a)]
  \item The sheaf $\eR^q f_! \sF$ is deduced by extension of scalars
    from $\dF_q$ to $\dF$ from the sheaf $\eR^q f_{0!} \sF_0$ on $Y_0$. As
    such, it is equipped with a Frobenius correspondence
    $F^\ast:F^\ast\eR^q f_! \sF \to \eR^q f_! \sF$. It is deduced by
    functoriality of $\eR f_!$ from the Frobenius correspondence of $\sF$:
    it is the composite 
    \[\xymatrix{
      F^\ast \eR^q f_! \sF \ar[r]^-\sim 
        & \eR^q f_! F^\ast \sF \ar[r]^-\sim 
        & \eR^q f_! \sF \text{.}
    }\]
  \item The trace formula for $\sF_q$ reads 
    \begin{align*}\tag{1}\label{VI:eq:1}
      \sum_{x\in X_0(\dF_q)} \tr(F_x^\ast,\sF_0) = \tr(F^\ast,\h^\bullet(X,\sF)) \text{;}
    \end{align*}
    that for $\eR^q f_! \sF$ reads 
    \begin{align*}\tag{2$_q$}\label{VI:eq:2_q} 
      \sum_{y\in Y_0(\dF_q)} \tr(F_y^\ast,\eR^q f_! \sF) = \tr(F^\ast, \h^\bullet(Y,\eR^q f_! \sF)) \text{.}
    \end{align*}
\end{enumerate}

We have, for $y\in Y(\dF)$, $(\eR^q f_! \sF)_y = \h_c^q(f^{-1}(y),\sF)$, and if
$y\in Y_0(\dF_q)$, 
\[
  \tr(F_y^\ast,\eR^q f_! \sF) = \tr(F^\ast,\h_c^q(f^{-1}(y),\sF)) \text{,}
\]
whence by the trace formula on the fibre 
\[
  \sum_{\substack{x\in X_0(\dF_q) \\ f(x) = y}} \tr(F_x^\ast,\sF_0) = \sum (-1)^q \tr(F_y^\ast,\eR^q f_! \sF) \text{.}
\]
Taking the alternating sum of the identities \eqref{VI:eq:2_q}, we therefore find 
\begin{align*}\tag{2}\label{VI:eq:2}
  \sum_{y\in Y_0(\dF_q)} \sum_{\substack{x \in X_0(\dF_q) \\ f(x) = y}} \tr(F_x^\ast,\sF_0) = \tr(F^\ast,\h^\bullet(Y,\eR^\bullet f_! \sF)) \text{.}
\end{align*}

Equality of the left-hand sides of \eqref{VI:eq:1} and \eqref{VI:eq:2}
follows from \eqref{VI:eq:2-3-1}, that of the right-hand sides from
\eqref{VI:eq:2-3-1*} (bearing in mind that $F^\ast$ is an endomorphism of
the whole spectral sequence). 





\subsection{}\label{VI:2-4}

Let $(A_i)_{i\in I}$ be a finite family of finite sets,
$A=\prod_{i\in I} A_i$, $\varepsilon_i$ a function on $A_i$, and
$\varepsilon$ the function $a\mapsto \prod \varepsilon_i(a_i)$ on $A$. We have 
\begin{align*}\tag{2.4.1}\label{VI:eq:2-4-1}
  \sum_{a\in A} \varepsilon(a) = \prod_{i\in I} \sum_{a_i\in A_i} \varepsilon_i(a_i) \text{.}
\end{align*}





\subsection*{2.4*}\label{VI:2-4*}

Let $(X_i)_{i\in I}$ be a finite family of separated schemes of finite type
over $k$ algebraically closed, $X=\prod_{i\in I} X_i$, $\sF_i$ an
$E_\lambda$-sheaf on $X_i$ and $\sF$ the external tensor product
$\bigboxtimes_{i\in I}\sF_i = \bigotimes_{i\in I} \pr_i^\ast(\sF_i)$
of the $\sF_i$. We have the K\"unneth formula 
\begin{align*}\tag{2.4.1*}\label{VI:eq:2-4-1*}
  \h_c^\bullet(X,\sF) = \bigotimes_{i\in I} \h_c^\bullet(X_i,\sF_i) \text{.}
\end{align*}

If one wants an isomorphism \eqref{VI:eq:2-4-1*} which is canonical, and independent of the
choice of a total order on $I$, one must take on the right-hand side the
graded tensor product in the sense of the Koszul rule
(\hyperref[IV]{Cycle}, \ref{V:1-3}). 





\subsection{}\label{VI:2-5}

Let $A$ be a finite set, $B$ a subset of $A$ and $\varepsilon$ a
function on $A$. We have 
\begin{align*}\tag{2.5.1}\label{VI:2-5-1}
  \sum_{a\in A} \varepsilon(a) = \sum_{a\in A\setminus B}\varepsilon(a) + \sum_{a\in B} \varepsilon(a) \text{.}
\end{align*}





\subsection*{2.5*}\label{VI:2-5*}

Let $X$ be a separated scheme of finite type over $k$ algebraically closed,
$Y$ a closed subscheme and $\sF$ an $E_\lambda$-sheaf on $X$. We have
a long exact cohomology sequence 
\begin{align*}\tag{2.5.1*}\label{VI:eq:2-5-1*}
\xymatrix{
  \cdots \ar[r]^-\partial 
    & \h_c^i(X\setminus Y,\sF) \ar[r] 
    & \h_c^i(X,\sF) \ar[r] 
    & \h_c^i(Y,\sF) \ar[r]^-\partial 
    & \cdots 
}
\end{align*}

More generally, starting from a finite filtration of $X$ by closed
subsets $X_p$ ($p\in \dZ$; we suppose that $X_p\supset X_{p+1}$, that $X_p=X$
for $p$ small enough, and that $X_p=\varnothing$ for $p$ large enough), we have a
spectral sequence 
\begin{align*}\tag{2.5.2*}\label{VI:eq:2-5-2*}
  E_1^{p q} = \h_c^{p+1}(X_p\setminus X_{p+1},\sF) \Rightarrow \h_c^{p+q}(X,\sF) \text{.}
\end{align*}





\subsection{}\label{VI:2-6}

Let $A$ be a finite set, $(A_i)_{i\in I}$ a finite covering of $A$ and
$\varepsilon$ a function on $A$. For $J\subset I$, let $A_J$ be
the intersection of the $A_j$ ($j\in J$). We have 
\begin{align*}\tag{2.6.1}\label{VI:eq:2-6-1}
  \sum_{J\subset I} (-1)^{|J|} \sum_{a\in A_J} \varepsilon(a) = 0 \text{.}
\end{align*}





\subsection*{2.6*}\label{VI:2-6*}

Let $X$ be a separated scheme of finite type over $k$ algebraically closed,
$(X_i)_{i\in I}$ a finite closed (resp. open) covering of $X$ by
subschemes, and $\sF$ an $E_\lambda$-sheaf on $X$. For $J\subset I$,
let $X_J$ be the intersection of the $X_j$ ($j\in J$). We have respective
(Leray) spectral sequences 
\begin{align*}\tag{2.6.1*}\label{VI:eq:2-6-1*}
  E_1^{p q} &= \bigoplus_{|J|=p+1>0} \h_c^q(X_J,\sF) \otimes_\dZ \textstyle\bigwedge^{|J|} \dZ^J \Rightarrow \h_c^{p+q}(X,\sF) \\ 
  \tag{2.6.2*}\label{VI:eq:2-6-2*}
  E_1^{p q} &= \bigoplus_{|J|=1-p>0} \h_c^q(X_J,\sF) \otimes_\dZ \textstyle\bigwedge^{|J|} \dZ^J \Rightarrow \h_c^{p+q}(X,\sF)
\end{align*}
If $J=\varnothing$ were not excluded, we would likewise have
spectral sequences converging to $0$. The factors $\bigwedge^{|J|} \dZ^J$ are there
to spare us from choosing a total order on $I$. 





\subsection{}\label{VI:2-7}

Let $A$ be a finite commutative group, and $\chi:A \to E_\lambda^\times$ a
non-trivial character. We have 
\begin{align*}\tag{2.7.1}\label{VI:eq:2-7-1}
  \sum_{a\in A} \chi(a) = 0 \text{.}
\end{align*}
We will prove the cohomological analogue of \eqref{VI:eq:2-7-1} by transposing
the following proof: if $x\in A$, $a\mapsto x a$ is a permutation of $A$, and 
\begin{align*}
  \sum_{a\in A}\chi(a) = \sum_{a\in A} \chi(x a) &= \chi(x) \sum_{a\in A} \chi(a) && \text{, whence} \\
  (\chi(x)-1)\sum_{a\in A} &= 0 \text{,}
\end{align*}
and there exists by hypothesis $x$ such that $\chi(x)-1\ne 0$. 





\begin{theorem*}[2.7*]\label{VI:2-7*}
Let $G_0$ be a connected commutative algebraic group over $\dF_q$ and
$\chi:G_0(\dF_q) \to E_\lambda^\times$ a non-trivial character. We have 
\begin{align*}\tag{2.7.1*}\label{VI:eq:2-7-1*}
  \h_c^\bullet(G,\sF(\chi)) = 0 \text{.}
\end{align*}
\end{theorem*}

For $x$ a rational point of $G$, denote by $t_x$ the translation $t_x(g) = x g$
of $G$. Formula $\fL t_x = t_{\fL(x)} \fL$ expresses that $(t_x,t_{\fL(x)})$
is an automorphism of the diagram $G\xrightarrow\fL G$ ($G$ equipped with the Lang
torsor). Let $\rho(g)$ be the automorphism of $(G,\sF(\chi))$ deduced from it.
For $g\in G_0(\dF_q)$, i.e. for $\fL(g) = e$, it is the identity on $G$, and
multiplication by $\chi(g)^{-1}$ on $\sF(\chi)$.

Denote by $\rho_H(g)$ the automorphism of $\h_c^\bullet(G,\sF(\chi))$ deduced from
$\rho(g)$. A homotopy argument (Lemma \ref{VI:2-8} below) shows that
$\rho_H(g) = \rho_H(e)$, hence is the identity. On the other hand, for
$g\in G_0(\dF_q)$, $\rho_H(g)$ is multiplication by
$(\chi^{-1}(g):\text{ on }\h^\bullet(G,\sF(\chi))$, multiplication by
$(\chi^{-1}(g)-1)$ is zero. Taking $g$ such that $\chi(g)\ne 1$, we obtain
\ref{VI:2-7*}. 





\begin{lemma_}\label{VI:2-8}
Let $X$ and $Y$ be two schemes over $k$ algebraically closed, with $X$
separated of finite type and $Y$ connected. Let $\sF$ be a sheaf on $X$ and
$(\rho,\varepsilon)$ a family of endomorphisms of $(X,\sF)$ parametrised
by $Y$:
\begin{align*}
  \rho:Y\times_k X &\to Y\times_k X && \text{is a $Y$-morphism, and} \\
  \varepsilon:\rho^\ast \pr_2^\ast \sF &\to \pr_2^\ast\sF && \text{a morphism of sheaves.}
\end{align*}
We suppose that $\rho$ is proper. For $y\in Y(k)$, let $\rho_H(y)^\ast$
be the endomorphism of $\h_c^\bullet(X,\sF)$ induced by $\rho_y:X\to X$ and
$\varepsilon_y:\rho_y^\ast\sF\to \sF$. Then $\rho_H(y)^\ast$ is
independent of $y$.
\end{lemma_}

Indeed, $\eR^p \pr_{1!} \pr_2^\ast\sF$ is the
constant sheaf on $Y$ with value $\h^p(X,\sF)$, and $\rho_H(y)^\ast$ is the
fibre at $y$ of the endomorphism 
\[\xymatrix{
  \eR^p \pr_{1!} \pr_2^\ast \sF \ar[r]^-{\rho^\ast} 
    & \eR^p \pr_{1!} \rho^\ast \pr_2^\ast \sF \ar[r]^-\varepsilon 
    & \eR^p \pr_{1!} \pr_2^\ast \sF 
}\]
of this sheaf.




\subsection{Remark}\label{VI:2-9}

Let $G_0$ be a connected algebraic group over $\dF_q$ and
$\rho:G_0(\dF_q) \to \operatorname{GL}(V)$ a linear representation
such that $V^{G_0(\dF_q)}=0$. We still have $\h_c^\bullet(G,\rho(L))=0$. 










\section{Sums in one variable}\label{VI:3}





\subsection{}\label{VI:3-1}

Weil was the first to have applied ``cohomological'' methods to
the study of exponential sums in one variable; since he had
proved the analogue of the Riemann hypothesis for Artin $L$-functions
over function fields, these methods provided him with excellent
estimates (in $O(\text{square root of the number of terms})$), and the
behaviour of these sums under ``extension of the base field.''

The bulk of this section is an exposition, in cohomological language, of
his results. 





\subsection{}\label{VI:3-2}

Cohomological methods lead here to computing groups
$\h_c^i(X,\sF)$, for $\sF_0$ an $E_\lambda$-sheaf on a curve $X-0$ over
$\dF_q$. These groups vanish for $i\ne 0,1,2$ and for
$i=0,2$, they have a simple interpretation (\ref{VI:1-8} a,b,c). Moreover, the
Euler-Poincaré characteristic
$\chi_c(\sF) = \sum (-1)^i \dim \h_c^i(X,\sF)$ (moreover equal to
$\chi(\sF) = \sum (-1)^i \dim \h^i(X,\sF)$ can be computed in terms of
$X$, of the rank of $\sF$ at the generic points of $X$, and of the
ramification properties of $\sF$. The essential result is the following.

Let $\bar X$ be a smooth projective connected curve, of genus $g$, over an
algebraically closed field $k$, $X$ a dense open of $\bar X$,
the complement of a finite set $S$ of points and $\sF$ a
smooth $E_\lambda$-sheaf on $X$. Let $\chi(X)=2-2 g-\# S$ be the
Euler-Poincaré characteristic of $X$, and $\operatorname{rg}(\sF)$ the
rank of $\sF$. For each point $s\in S$, we define an integer
$\swan_s(\sF)$, the Swan conductor, measuring the wild ramification of
$\sF$, and 
\begin{align*}\tag{3.2.1}\label{VI:eq:3-2-1}
  \chi_c(\sF) = \operatorname{rg}(\sF) \cdot \chi(X) - \sum_{s\in S} \swan_s(\sF) 
\end{align*}
\cite{ra65}. 





\subsection{}\label{VI:3-3}

Let $\bar X$, $j:X\hookrightarrow\bar X$ and $\sF$ be as in \ref{VI:3-2}. The
Poincaré duality theorem takes the very convenient form:
$\h^i(\bar X,j_\ast \sF)$ and $\h^{2-i}(\bar X,j_\ast(\sF^\vee(1)))$ are
dual to each other (\hyperref[V]{Dualité}, \ref{V:1-3}). 





\subsection{}\label{VI:3-4}

Put $k=\dF$, and suppose that $X$, $\bar X$, and $\sF$ come from
$X_0$, $\bar X_0$ and $\sF_0$ over $\dF_q$. If $\sF_0$ becomes trivial on a
finite covering of $X_0$, Weil proved that the complex conjugates of the
eigenvalues of $F^\ast$ on $\h^i(\bar X,j_\ast\sF)$ are of absolute
value $q^{i/2}$. Such sheaves $\sF_0$ correspond to Artin
$L$-functions over the function field of $\bar X_0$. 





\subsection*{3.5}\label{VI:3-5_} % originally one of two 3.5s

Under the same hypotheses, or more generally if $E_\lambda$ is the
completion at a place $\lambda$ of a number field $E$ and $\sF_0$
belongs to an infinite compatible system of $v$-adic representations
($v$ a place of $E$), one can compute 
\[
  \prod_i \det(-F^\ast,\h^i(X,\sF))^{(-1)^i} 
\]
from \emph{local} information on $\sF_0$. This is explained in
\cite{de73}. For $\sF_0$ trivial on a finite covering of $X_0$, this is
the expression of the constant in the functional equation of an Artin
$L$-function as a product of local constants. 





\subsection{Example}\label{VI:3-5}

Let $\psi$ be the character
$\exp\left(\frac{2\pi i}{p} \tr_{\dF_q/\dF_p}\right)$ of $\dF_q$; we have
$\psi(a^p-a)=0$. Let $X_0$ be a smooth projective absolutely
irreducible curve of genus $g$ over $\dF_q$, and $f$ a rational function on
$X_0$, i.e. a morphism $f:X_0 \to \dP^1$, not identically equal to
$\infty$. We are interested in the sum 
\[
  S_f' = \sum_{\substack{x\in X_0(\dF_q) \\ f(x)\ne \infty}} \psi(f(x)) 
  \text{.}
\]
This sum is zero for $f$ of the form $g^p-g$. This suggests
modifying the sum $S_f'$ as follows: 
\begin{enumerate}[\indent a)]
  \item for every closed point $x$ of $X_0$, we set $v_x(f) = $ order of the
    pole of $f$ at $x$ if $f(x)=\infty$, $v_x(f)=0$ otherwise, and
    $v_x^\ast(f) = \inf v_x(f+g^p-g)$ (infimum over $g$).
  \item If $v_x^\ast(f)=0$, $x\in X_0(\dF_q)$, and $f+g^p-g$ is
    regular at $x$, we set $\psi(f(x)) = \psi((f+g^p-g)(x))$. We finally set 
    \begin{equation*}\tag{3.5.1}\label{VI:eq:3-5-1}
      S_f = \sum'_{x\in X_0(\dF_q)} \psi(f(x)) \text{,}
    \end{equation*}
    where $(-)'$ indicates that the sum is extended over those $x$ with
    $v_x^\ast(f)=0$.
\end{enumerate}

We have $S_f=S_{f+g^p-g}$. We propose to check that if $f$ is not of the
form $g^p-g+c^{t e}$, then 
\begin{equation*}\tag{3.5.2}\label{VI:eq:3-5-2}
  |S_f|\leqslant \left(2 g-2+\sum_{v_x^\ast(f)\ne 0} [k(x):\dF_q] (1+v_x^\ast(f))\right) q^{1/2} \text{.}
\end{equation*}

We begin by passing from the complex to the $\ell$-adic, replacing $\psi$ by
a character of the form $\psi_0\circ \tr_{\dF_q/\dF_p}$, with
$\psi_0:\dF_p\hookrightarrow E_\lambda^\times$. For $j:U_0\hookrightarrow X_0$
the inclusion of the open where $f\ne\infty$, we then prove that 
\begin{equation*}\tag{3.5.3}\label{VI:eq:3-5-3}
  S_f = \sum (-1)^i \tr(F^\ast,\h^i(X,j_\ast \sF(\psi f))) \text{.} 
\end{equation*}

The trace formula reduces this statement to the following (for
$x\in X_0(\dF_q)$): 
\begin{enumerate}[\indent a)]
  \item If $v_x^\ast(f)=0$, then $F_x^\ast$ on $j_\ast \sF(\psi f)$ is
    $\psi(f(x))$;
  \item If $v_x^\ast(f)\ne 0$, then $j_\ast \sF(\psi f)$ ramifies at $x$:
    $(j_\ast \sF(\psi f))_{\bar x}=0$.
\end{enumerate}

Formula a) follows from \eqref{VI:eq:1-7-6} if $v_x(f)=0$, and we reduce
to this case replacing $f$ by $f+g^p-g$: up to isomorphism, this does not
change $\sF(\psi f)_0$ on the open where $f$ and $g$ are regular
(\ref{VI:1-3}).

Formula b) follows from 
\begin{equation*}\tag{3.5.4}\label{VI:eq:3-5-4}
  \swan_x \sF(f,\psi) = v_x^\ast(f) \text{.} 
\end{equation*}
After reduction to the case where $v_x(f)=v_x^\ast(f)$ (whence $p\nmid v_x(f)$)
and extension of scalars to $\dF$, this formula is in
\cite[4.4]{se61}.

Finally, we check that the sheaf $\sF(\psi f)$ on $U$ is non-constant if
$f$ is not of the form $g^p-g+c^{t e}$: since $\psi_0$ is injective, the
triviality of $\sF(\psi f)=\sF(\psi_0 f)$ (\ref{VI:1-8}.ii) is equivalent to
that of the $\dF_q$-torsor of equation would be the inverse image of a
$\dF_q$-torsor on $\spec(\dF_q)$, of equation $T^p-T-\lambda=0$, the
$\dF_p$-torsor of equation $T^p-T-(f-\lambda)=0$ on $U_0$ would be trivial
and $f$ would therefore be of the form $g^p-g+\lambda$. The
$\h^i(\bar X,j_\ast \sF(\psi f))$ therefore vanish for $i\ne 1$, formula
\eqref{VI:eq:3-5-3} reduces to 
\begin{equation*}\tag{3.5.5}\label{VI:eq:3-5-5}
  S_f = -\tr(F^\ast,\h^1(\bar X,j_\ast \sF(\psi f))) 
\end{equation*}
and, by \eqref{VI:eq:3-2-1} and \eqref{VI:eq:3-5-4}, the $\h^1$ is of
dimension $2 g-2+\sum_{v_x^\ast(f)>0} [k(x):\dF_q](1+v_x^\ast(f))$ and $-S_f$
is the sum of this number of eigenvalues of $F^\ast$, each of complex
absolute value $q^{1/2}$. 





\subsection{}\label{VI:3-6}

Suppose that, for an automorphism $\sigma$ of $X_0$, we have
\[
  f(\sigma x) = - f(x) \text{.}
\]
The sum $S_f$ is then real. If $\sigma$ is involutive, Poincaré
duality makes it possible to say a little more: putting
$\sF_0=j_\ast \sF(\psi f)$, $\sG_0 = j_\ast \sF(\psi(-f))$, we have
\begin{enumerate}[\indent a)]
  \item $\sF_0$ and $\sG_0$ are in duality,
  \item $\sigma^\ast \sF_0 \iso \sG_0$ and $\sigma^\ast \sG_0 \iso \sF_0$, for
    natural isomorphisms such that the composite
    $\sF_0=(\sigma^2)^\ast \sF_0 \iso \sigma^\ast \sG_0 \iso \sF_0$ is
    the identity.
  \item The pairing $\sF_0\otimes \sG_0 \to E_\lambda$ satisfies
    $\sigma^\ast(f\cdot g) = \sigma^\ast(f) \cdot \sigma^\ast (g)$.
\end{enumerate}

Passing to cohomology, we find that $\h^1(X,\sF)$ and $\h^1(X,\sF)$ are in
perfect duality with values in $E_\lambda(1)$, and that the bilinear form
$\alpha\cdot \sigma^\ast\beta$ on $\h^1(X,\sF)$ is \emph{alternating}:
$\sigma^\ast$ acts trivially on $E_\lambda(1)$, and 
\[
  \alpha\cdot \sigma^\ast \beta = \sigma^\ast(\alpha\cdot \sigma^\ast \beta) = \sigma^\ast \alpha\cdot \beta = -\beta\cdot \sigma^\ast \alpha \text{.} 
\]
We conclude that $\h^1(X,\sF)$ is of even dimension (one checks moreover
easily that each non-zero $v_x^\ast(f)$ is odd) and that the eigenvalues
of $F^\ast$ on $\h^1(X,\sF)$ are grouped in pairs $\alpha$ and
$q/\alpha$. 





\subsection{Example}\label{VI:3-7}

This applies to Kloosterman sums
$\sum_{x\in \dF_q^\times} \psi\left(x+\frac a x\right)$ (take
$X_0=\dP^1$, $f=x+\frac a x$, $\sigma x=-x$; the poles of $f$ are $0$ and
$\infty$, and at each of them $v_x^\ast(f)=1$; the $\h^1$ is of dimension
$-2+2+2=2$). We therefore have (for $a\ne 0$) 
\[
  \sum_{\substack{x y=a \\ x,y\in \dF_{q^n}}} \exp\left(\frac{2\pi i}{p} \tr_{\dF_{q^n}/\dF_p}(x+y)\right) = (-\alpha^n+\alpha^{-n}) \qquad \text{, }\alpha\bar \alpha = q \text{.} 
\]
This result is due to L.\ Carlitz \cite{ca69}.

Here is now an application of a Lang-Weil method. 





\begin{proposition_}\label{VI:3-8}
Let $P\in \dF_q[X_1,\dots,X_n]$ be a polynomial in $n$ variables, of degree
$d$, which we suppose is not of the form $Q^p-Q+C^{te}$. Still putting
$\psi(x) = \exp\left(\frac{2\pi i}{p} \tr_{\dF_q/\dF_p}(x)\right)$ we have
\[
  \left|\sum_{x_i\in \dF_q} \psi P(x_1,\dots,x_n)\right| \leqslant (d+1) q^{n-\frac 1 2} \text{.}
\]
\end{proposition_}

For the left-hand side we have the trivial estimate $q^n$. It therefore suffices to
prove \ref{VI:3-8} when $d-1<q^{1/2}$; let us only suppose that
$d<q+1$, and let $P_d$ be the homogeneous part of degree $d$ of $P$. Recall the 





\begin{lemma_}\label{VI:3-9}
A hypersurface of degree $d$ in $\dP^r$ ($r\geqslant i$) cannot pass
through all rational points over $\dF_q$ unless $d\geqslant q+1$.
\end{lemma_}

We argue by induction: if there exists a rational hyperplane not
entirely contained in the hypersurface (such is not the case for $r=1$), we
apply the induction hypothesis to the trace of the hypersurface on this
hyperplane. Otherwise, the degree $d$ is $\geqslant$ the number of rational
hyperplanes, $\geqslant q+1$.

Let us now distinguish two cases.


\paragraph{Case 1: $p\nmid d$.}
Applying the lemma to $P_d$, we see that, up to a linear
change of variables, we may suppose that $P_d(1,0,\ldots,0)\ne 0$, i.e.
that the coefficient of $X_1^d$ in $P$ is non-zero. A polynomial in $1$
variable 
\[
  S(X) = \sum_{i=0}^d a_i x^i \text{,} 
\]
with $a_d\ne 0$ ($p\nmid d$) is never of the form $Q^p-Q+C^{t e}$, and
the estimate \eqref{VI:eq:3-5-2} reduces to
\[
  \left| \sum \psi(S(x)) \right| \leqslant (d-1) q^{1/2} \text{.}
\]
Applying this estimate to the partial sums obtained by varying only
$x_1$, we find 
\[
  \left| \sum \psi P(x_1,\dots,x_n) \right| = \left| \sum_{x_2,\dots,x_n} \sum_{x_1} \psi P(x_1,\dots,x_n)\right| \leqslant q^{n-1} (d-1) q^{1/2} 
\]
as promised.


\paragraph{Case 2: $p\mid d$ ($d>0$).}
If $P_d$ is a $p$-th power, replacing $P$ by $P-(P_d-P_d^{1/p})$ does not
change the sum considered, and lowers the degree: we dispose of
this case by induction on $d$. Otherwise, the differential of
$P_d$ is not identically zero; applying \ref{VI:3-9}, we may suppose,
up to a linear change of variables, that it is non-zero
at the point $(1,0,\dots,0)$. If we write 
\[
  P_d = \sum_{i=0}^d X_1^{d-i} S_i(X_2,\dots,X_n) \text{,} 
\]
this means that the linear form $S_1$ is not identically zero.

The form $S_0$ is a constant, and replacing $P$ by
$P-(S_0 X_1^d -S_0^{1/p} X_1^{d/p})$, we may suppose that it is zero. Let
finally $-\lambda$ be the coefficient of $X_1^{d-1}$ in $P$. For fixed
$x_2,\dots,x_n$, we have 
\[
  P(X_1,x_2,\dots,x_n) = (S_1(x_2,\dots,x_n)-\lambda) X_1^{d-1} + \text{lower-degree terms in $X_1$}
\]
and if $S_1(x_2,\dots,x_n)\ne \lambda$, we therefore have 
\[
  \left| \sum_{x_1} \psi P(x_1,\dots,x_n)\right| \leqslant (d-2) q^{1/2} \text{.} 
\]
In total, 
\begin{align*}
  \left|\sum \psi P(x_1,\dots,x_n)\right| 
    &\leqslant \left|\sum_{S=\lambda} \psi P(x_1,\dots,x_n)\right| + \sum_{S(X_2,\dots,X_n)\ne\lambda} \left| \sum_{x_1} \psi P(x_1,\dots,x_n)\right| \\
    &\leqslant q^{n-1} + (q^{n-1}-q^{n-2})(d-2) q^{1/2} \\
    &< (d-2) q^{n-1/2} + q^{n-1} \\
    &< (d-1) q^{n-1/2} \text{.}
\end{align*}

Another result of this nature is given by R.A.\ Smith
\cite{sm70}. 










\section{Gauss sums and Jacobi sums}\label{VI:4}





\subsection{}\label{VI:4-1}

Let $k$ be a finite field of characteristic $p$, $\chi$ a character of
$k^\times$, and $\psi$ a non-trivial character of the additive group of $k$. We
take as definition of Gauss sums: 
\begin{equation*}\tag{4.1.1}\label{VI:eq:4-1-1}
  \tau(\chi,\psi) = -\sum_{x\in k^\times} \psi(x) \chi^{-1}(x) 
\end{equation*}
(note the sign). Classically, $\chi$ and $\psi$ take complex values.
We will take them with values in $E_\lambda^\times$ (cf. the proof of
\ref{VI:1-15}). Let us regard $-\tau(\chi,\psi)$ as a sum over the
rational points of the scheme $\dG_m$ over $k$. By section \ref{VI:1},
we have $-\tau(\chi,\psi) = \tr(F^\ast,\h_c^\ast(\dG_m,\sF(\psi \chi^{-1}))$. Moreover 





\begin{proposition_}\label{VI:4-2}
The cohomology of $\dG_m$ with coefficients in $\sF(\psi\chi^{-1})$ satisfies
\begin{enumerate}[(i)]
  \item $\h_c^i=0$ for $i\ne 1$, and $\dim \h_c^1=1$.
  \item $F^\ast$, acting on $\h_c^1$, is multiplication by
    $\tau(\chi,\psi)$.
  \item If $\chi$ is non-trivial, we have $\h_c^\bullet \iso \h^\bullet$.
\end{enumerate}
\end{proposition_}
\begin{proof}[Proof]
For every $n$ prime to $p$, denote by $K_n$ the $\dmu_n$-torsor on $\dG_m$
defined by the K\"ummer exact sequence
$0 \to \dmu_n \to \dG_m \xrightarrow{x^n} \dG_m \to 0$. If $k$ has $q$
elements, we have over $k$: $\dmu_{q-1}=k^\times$, and the Lang torsor on
$\dG_m/k$ is $K_{q-1}$. Hence $\sF(\chi^{-1})=\chi(K_{q-1})$.
Assertion \ref{VI:4-2}(ii) follows from (i,iii), themselves contained in
the following geometric statement, where $\sF(\psi)$ denotes the sheaf
on $\dG_m/\dF$ deduced by extension of scalars from $k$ to $\dF$ of
$\sF(\psi)$ on $\dG_m/k$ (\ref{VI:1-8}.b).
\end{proof}





\begin{proposition_}\label{VI:4-3}
Let $\chi:\dmu_n \to E_\lambda^\times$. The cohomology of $\dG_m$ with
coefficients in $\sF(\psi)\otimes \chi(K_n)$ satisfies
\begin{enumerate}[\indent (i)]
  \item $\h_c^i=0$ for $i\ne 1$, and $\dim \h_c^1{} = 1$.
  \item If $\chi$ is non-trivial, we have $\h_c^\bullet{}\iso \h^\bullet{}$.
\end{enumerate}
\end{proposition_}
\begin{proof}[Proof]
The sheaf $\sF(\psi)$ is the restriction to $\dG_m$ of a locally
constant sheaf on $\dG_a$, wildly ramified at infinity, of Swan
conductor $1$. The sheaf $\chi(K_n)$ is constant if $\chi=1$; if $\chi\ne 1$, it
is ramified at $0$ and $\infty$, tamely.

The sheaf $\sF(\psi)\otimes\chi(K_n)$ is therefore ramified at $\infty$;
applying \ref{VI:1-18}(b,c), we find that
$\h_c^i(\dG_m,\sF(\psi)\otimes \chi(K_n))=0$ for $i\ne 1$. If $\chi\ne 1$, it
is ramified at $0$ and $\infty$ and (ii) follows from \ref{VI:1-19}.a.

The Swan conductors are $0$ at $0$ and $1$ at $\infty$. By
\eqref{VI:eq:3-2-1}, the Euler-Poincaré characteristic is therefore $-1$ and
this completes the proof.
\end{proof}





\subsection{Remark}\label{VI:4-4}

If $\chi$ is non-trivial, \ref{VI:4-2}(iii) and Poincaré duality
show that $\h_c^1(\dG_m,\sF(\psi\chi^{-1}))$ and
$\h_c^1(\dG_m,\sF(\psi^{-1}\chi))$ are in duality (duality with values
in $E_\lambda(-1)$). We therefore have
$\tau(\chi,\psi)\cdot \tau(\chi^{-1},\psi^{-1}) = q$, i.e.
$|\tau(\chi,\psi)|=1$. 





\subsection{}\label{VI:4-5}

If $k$ is an extension of degree $N$ of $\dF_q$, one can also regard
\eqref{VI:eq:4-1-1} as a sum in $N$ variables over $\dF_q$. More
generally, let $k$ be an étale algebra over $\dF_q$, of degree $N$ over
$\dF_q$. It is a product of fields $k_i$, of degree $N_i$, and we set 
\begin{equation*}\tag{4.5.1}\label{VI:eq:4-5-1}
  \varepsilon(k) = (-1)^{\sum (N_i+1)} \text{.} 
\end{equation*}
It is the signature of the permutation of $S=\hom_{\dF_q}(k,\dF)$ induced by the
Frobenius substitution $\varphi\in \gal(\dF/\dF_q)$.

Let $\psi:\dF_q\to E_\lambda^\times$ be a non-trivial character, and
$\chi:k^\times \to E_\lambda^\times$. We set 
\begin{equation*}\tag{4.5.2}\label{VI:eq:4-5-2}
  \tau_{\dF_q}(\chi,\psi) = (-1)^N \sum_{x\in k^\times} \psi \tr_{k/\dF_q}(x) \cdot \chi^{-1}(x) \text{.} 
\end{equation*}

If $\chi$ has coordinates the $\chi_i:k_i^\times \to E_\lambda^\times$, we
have the trivial identity 
\begin{equation*}\tag{4.5.3}\label{VI:eq:4-5-3}
  \tau_{\dF_q}(\chi,\psi) = \varepsilon(k) \prod \tau(\chi_i,\psi\circ \tr_{k_i/\dF_q}) \text{.}
\end{equation*}

After some preliminaries, we will give in \ref{VI:4-10} a
cohomological interpretation of the sums \eqref{VI:eq:4-5-2}. In
\ref{VI:4-12}, we will interpret the Hasse-Davenport identity as a
twisted form (cf. \ref{VI:1-12}) of the following special case of
\eqref{VI:eq:4-5-3}: for $k=\dF_q^N$, and $\chi$ a character of
$\dF_q^\times$, we have
$\tau_{\dF_q}(\chi\circ N_{k/\dF_q},\psi)=\tau(\chi,\psi)^N$. 





\subsection{}\label{VI:4-6}

Recall that for $M$ a finite-type projective module over a ring $A$, the
functor $\spec(B)\mapsto M\otimes_A B$, from affine schemes over $\spec(A)$
to $\mathsf{Ens}$ is represented by the affine scheme $\dV(M^\vee)$
(notations of EGA), the spectrum of $\operatorname{Sym}_A^\bullet(M^\vee)$. If
$M=A^n$, it is the typical affine space of dimension $n$.

Suppose that $M$ is a unital $A$-algebra, and put
$V=\dV(M^\vee)$. By definition, for every extension $B$ of $A$, the set
$V(B)$ of points of $V$ with coordinates in $B$ is $M\otimes_A B$. The
functor $B\mapsto M\otimes_A B$ taking values in unital
rings, $V$ is a ring scheme with units over $\spec(A)$. The
norm and trace morphisms: $M\otimes_A B\to B$ are functorial in $B$; they
therefore correspond to scheme morphisms $N$ and $T$ from $V$ to
$\dG_a$. We will have to consider the following schemes deduced from $V$:
\begin{enumerate}[a)]
  \item $V^\ast$ is the open of invertible elements of $V$ (a group
    scheme for $\cdot$);
  \item $W$ is the hyperplane of equation $T=0$ and $W^\ast=W\cap V^\ast$;
  \item $P$ is the hyperplane at infinity $V\setminus \{0\}/\dG_m$ of the
    affine space $V$ over $\spec(A)$. If $Q$ is the hyperplane at infinity of $W$,
    $W^\ast/\dG_m$ and $V^\ast/\dG_m$ are opens of the projective spaces
    $Q$ and $P$ over $\spec(A)$. 
    \begin{equation*}\tag{4.6.1}\label{VI:eq:4-6-1}
    \xymatrix{
      W^\ast \ar@{^{(}->}[r] \ar[d]^-\pi 
        & V^\ast \ar[d]^-\pi \\
      W^\ast/\dG_m \ar@{^{(}->}[r] 
        & V^\ast/\dG_m 
    }\qquad\subset\qquad
    \xymatrix{
      W\setminus \{0\} \ar@{^{(}->}[r] \ar[d]^-\pi 
        & V\setminus \{0\} \ar[d]^-\pi \\
      Q \ar@{^{(}->}[r] 
        & P \text{.}
    }
    \end{equation*}
\end{enumerate}

If $M=A^I$, $I$ a finite set, then $V\simeq \dG_a^I$, $V^\ast=\dG_m^I$,
$V^\ast/\dG_m$ is the torus $\dG_m^I/(\dG_m\text{ diagonal})$, $N$ and $T$
being written $\prod x_i$ and $\sum x_i$, and $Q$ is therefore the
projective hyperplane of equation $\sum x_i=0$ of $P$.

The case of interest to us is that where $M$ is finite étale over $A$. The
preceding description then holds locally on $\spec(A)$ (for the
étale topology) and one can write $V^\ast=\dG_m^I$, for $I$ a
locally constant sheaf of finite sets on $\spec(A)$. If for example $A$ is a
field, of algebraic closure $\bar A$, and $M$ is a separable
$A$-algebra, we put $I=\hom_A(M,\bar A)$, and, over $\bar A$, we have
$V\sim \dG_a^I$. Via this isomorphism, $N$ and $T$ are written $\prod x_i$ and
$\sum x_i$. 





\subsection{Kummer torsor}\label{VI:4-7}

Let $S$ be a scheme, $n$ an integer invertible on $S$ and $G$ a commutative
group scheme with connected fibres (denoted multiplicatively) over $S$.
The sequence 
\[\xymatrix{
  0 \ar[r] 
    & G_n \ar[r] 
    & G \ar[r]^-{x^n} 
    & G \ar[r] 
    & 0 
}\]
is exact. It defines a $G_n$-torsor $K_n(G)$ (or simply $K_n$) on
$G$. For $\chi:G_n \to E_\lambda^\times$ a homomorphism from the étale
$S$-sheaf $G_n$ to the constant sheaf $E_\lambda^\times$, we denote by
$\sK_n(\chi)$ the $E_\lambda$-sheaf $\chi^{-1}(K_n)$.

The torsors $K_n$ form a projective system of torsors under the
projective system of groups $G_n$ (transition morphisms
$x\mapsto x^d:G_{n d}\to G_n$). This expresses the commutativity of the diagrams 
\[\xymatrix{
  0 \ar[r] 
    & G_{n d} \ar[r] \ar[d]^-{x^d} 
    & G \ar[r]^-{x^{n d}} \ar[d]^-{x^d} 
    & G \ar[r] \ar@{=}[d] 
    & 0 \\
  0 \ar[r] 
    & G_n \ar[r] 
    & G \ar[r]^-{x^n} 
    & G \ar[r] 
    & 0 \text{.} 
}\]
We therefore have $\sK_n(\chi) = \sK_{n d}(\chi\circ x^d)$. 





\subsection{}\label{VI:4-8}

Let us apply construction \ref{VI:4-6} with $A=\dF_q$, and $M=k$ an étale
algebra over $\dF_q$. In accordance with the general conventions, we will denote
with a subscript $0$ the $\dF_q$-schemes denoted $V,V^\ast,\ldots$ in
\ref{VI:4-6}. The same letters without subscript denote the schemes over
$\dF$ deduced from them by extension of scalars.

We put $I=\hom(k,\dF)$, whence $V\sim \dG_a^I$ and $T$ is written
$(x_i)\mapsto \sum x_i$. 
\begin{equation*}\tag{4.8.1}\label{VI:eq:4-8-1}
  \sF(\psi\circ \tr_{k/\dF_q}) = T^\ast \sF(\psi) = \bigotimes_i \pr_i^\ast \sF(\psi) \text{.} 
\end{equation*}
Since $V^\ast\sim \dG+m^I$, a character $\chi$ of $V_n^\ast\sim \dmu_n^I$,
with values in $E_\lambda^\times$, is written as a family
$(\chi_i)_{i\in I}$ of characters of $\dmu_n$ indexed by $I$. It is
defined over $\dF_q$ if, for every $\sigma\in \gal(\dF/\dF_q)$, we have
$\chi_{\sigma_i} = \chi_i\circ \sigma^{-1}$. It then defines an
$E_\lambda$-sheaf $\sK_n((\chi_i)_{i\in I})$ on $V_0^\ast$. If the product
of the $\chi_i$ is trivial, $\chi$ factors through a character of
$(V_0^\ast/\dG_m)_n$ and $\sK_n((\chi_i)_{i\in I})$ is the inverse image
of a sheaf, denoted likewise, on $V-0^\ast/\dG_m$. This construction
compares as follows with the Lang torsor. 





\begin{lemma_}\label{VI:4-9}
\begin{enumerate}[(i)]
  \item For $n$ sufficiently divisible, we have a commutative diagram
    \[\xymatrix{
      0 \ar[r] 
        & (V_0^\ast)_n \ar[r] \ar[d]^-\tau 
        & V_0^\ast \ar[r]^-{x^n} \ar[d]^-\tau 
        & V_0^\ast \ar[r] \ar@{=}[d] 
        & 0 \\
      0 \ar[r] 
        & k^\times \ar[r] 
        & V_0^\ast \ar[r]^-\fL 
        & V_0^\ast \ar[r] 
        & 0 \text{;} 
    }\]
    if $\chi$ is a character $\chi:k^\times \to E_\lambda^\times$, we have
    $\sF(\chi) = \sK_n(\chi\circ\tau)$.
  \item For $n$ sufficiently divisible, $\tau$ identifies $k^\times$ with the coinvariants
    of $\gal(\bar\dF/F)$ acting on $V_n^\ast\sim \dmu_n^I$.
  \item For $k$ a field, $N=[k:\dF_q]$, $\omega\in I$ an embedding of $k$
    into $\dF$, and $n=q^N-1$, the component of index $\omega$ of
    $\chi\circ\tau$ is $\chi\circ \omega^{-1}$.
  \item If $\chi$ is non-trivial on each factor of $k$, the
    $(\chi\circ \tau)_i$ are all non-trivial.
\end{enumerate}
\end{lemma_}

It suffices to prove the lemma when $k$ is a field, of degree $N$ over
$\dF_q$. In this case, $n$ is ``sufficiently divisible'' if $q^N-1\mid n$.
Choose an embedding $\omega$ of $k$ into $\dF_q$; $I$ then identifies
with $\dZ/N$: to $i\in \dZ/N$ corresponds $\omega_i=\omega^{q^i}$. Via
the isomorphism $V_0^\ast(\dF) = {F^\times}^I$, we have 
\begin{enumerate}[\indent a)]
  \item $x\in k^\times$ corresponds to $(\omega_i(x))\in {\dF^\times}^I$;
  \item $F((x_i)_{i\in \dZ/N})=(x_{i-1}^q)_{i\in \dZ/N}$.
\end{enumerate}

A character $\chi=(\chi_i)$ of $V_n^\ast\sim \dmu_n^I$ will be defined over
$\dF_q$ if $\chi_{-i} =\chi_0(x^{q^i})$ ($i\in \dZ$). If $q^N-1\mid n$, there are
$q^N-1$ such characters: $\chi_0$ factors through $\dmu_{q^N-1}$, and
determines the $\chi_i$. To prove (ii), it therefore suffices to check (i)
and (iii) (or its corollary (iv)) which ensures that $\chi\mapsto \chi\circ \tau$
is injective.

Let us prove (i) and (iii), for $n=q^N-1$. Let us denote additively the group of
endomorphisms of $V^\ast\sim \dG_m^I$; in particular, denote by $n$
the operator $x\mapsto x^n$. If $\alpha$ is the circular permutation
operator $(x_i)\mapsto (x_{i-1})$, we have
$q^N-1=(q\alpha)^N-1 = (q\alpha-1)((q\alpha)^{N-1}+\cdots + 1)$. This
determines $\tau$: we have $\tau((x_i)) = \prod_{0\leqslant j<N} x_{i-j}^{q^j}$,
and the map induced from $(V_0^\ast)_n$ to $k^\times$ is written [a)
above] $(x_i)\mapsto \prod x_{-i}^{q^i}$; (iii) follows. 





\begin{proposition_}\label{VI:4-10}
The cohomology of $V^\ast$ with coefficients in
$\sF(\chi^{-1}\cdot \psi\tr_{k/\dF_q})$ satisfies
\begin{enumerate}[\indent (i)]
  \item $\h_c^i = 0$ for $i\ne N$, and $\dim \h_c^N=1$.
  \item On $\h_c^N$, $F^\ast$ is multiplication by
    $\tau_{\dF_q}(\chi,\psi)$.
  \item If $\chi$ is non-trivial on each factor of $k$, we have
    $\h_c^\bullet \iso \h^\bullet$.
\end{enumerate}
\end{proposition_}

Let us apply \ref{VI:4-9}(i): over $\dF$, if $\chi\circ \tau=(\chi_i)_{i\in I}$,
we have $\sF(\chi) = \sK_n((\chi_i)_{i\in I})$. Points (i) and (iii)
therefore follow from the following more geometric statement (for (iii),
apply \ref{VI:4-9}(iv)) and (ii) follows from it by the trace formula: 





\begin{proposition_}\label{VI:4-11}
Let $(\chi_i)_{i\in I}$ be a family of characters of $\dmu_n(\dF)$. The
cohomology of $V^\ast\simeq \dG_m^I$ with coefficients in
$\sK_n((\chi_i)_{i\in I})\otimes \sF(\psi\tr_{k/\dF_q})$ satisfies
\begin{enumerate}[\indent (i)]
  \item $\h_c^i{}=0$ for $i\ne N$, and $\dim \h_c^N{} = 1$.
  \item If the $\chi_i$ are all non-trivial, we have
    $\h_c^\bullet{}\iso \h^\bullet$.
\end{enumerate}
\end{proposition_}

We have $V^\ast\sim \dG_m^I$,
$\sK((\chi_i)) = \bigotimes_i \pr_i^\ast \chi_i(K_n(\dG_m))$ and
$\sF(\psi\circ \tr_{\dF_q/k})=T^\ast \sF(\psi) = \bigotimes_i \pr_i^\ast \sF(\psi)$.
This makes it possible to apply the K\"unneth formula, and \ref{VI:4-11} follows from
\ref{VI:4-3}. 





\subsection{}\label{VI:4-12}

From \ref{VI:2-4*}* and (\hyperref[VI]{Cycle}, \ref{VI:1-3} example 2), we also
deduce that the group of permutations $\sigma$ of $I$ such that
$\chi_i = \chi_{\sigma i}$ ($i\in I$) acts on $\h_c^N$ by multiplication by
the signature $\varepsilon(\sigma)$. We are going to deduce from this a second proof
of the Hasse-Davenport identity.

If $\chi$ is a character of $\dF_q^\times$, the fact that, over $\dF$, $N$
is written $(x_i)\mapsto \prod x_i$ gives: 
$\sF(\chi\circ N)=N^\ast\sF(\chi) = \bigotimes_i \pr_i^\ast\sF(\psi)$ 
and the K\"unneth formula gives: 
\[
  \h_c^\bullet(V^\ast,\sF(\chi^{-1}\circ N,\psi\circ \tr)) \sim \h_c^\bullet(\dG_m,\sF(\chi^{-1}\psi))^{\otimes I} \text{,} 
\]
where, on the right-hand side, the tensor product is taken in the sense
\ref{VI:2-4*}*. Since $\h_c^1(\dG_m,\sF(\psi\chi^{-1}))$ is of dimension
$1$, its tensor power $\otimes I$, in the ordinary sense, depends only
on the cardinal $n$ of $I$. Applying (\hyperref[IV]{Cycle}, \ref{IV:1-3} example
2), we obtain a canonical isomorphism 
\begin{equation*}\tag{4.12.1}\label{VI:eq:4-12-1}
  \h_c^N\left(V,\sF(\chi^{-1}\circ N)\otimes T^\ast \sF(\psi)\right) \sim \h_c^1(\dG_m,\sF(\psi\chi^{-1})) \otimes_\dZ \textstyle \bigwedge^N \dZ^I \text{.}
\end{equation*}
To compute the action of $F^\ast$, the most convenient is to adopt the
Galois point of view and to say that the isomorphism \eqref{VI:eq:4-12-1} being
canonical, it is compatible with the action by transport of structure of
$\gal(\dF/\dF_q)$. If $\varepsilon(k)$ is the signature of the permutation
$\varphi$ of $I$, we find that 
\begin{align*} 
  \tau_{\dF_q}(\chi\circ N_{k/\dF_q},\psi) 
    &= \tr\left(\varphi^{-1},\h_c^N\left(V^\ast,N^\ast\sF(\chi^{-1})\otimes T^\ast \sF(\psi)\right)\right) \\ 
    &= \varepsilon(k) \tr\left(\varphi^{-1},\h_c^1(\dG_m,\sF(\chi^{-1}\psi))\right)^N \\
    &= \varepsilon(k) \tau(\chi,\psi)^N \text{,}
\end{align*}
i.e.
\begin{equation*}\tag{4.12.2}\label{VI:eq:4-12-2}
  \tau_{\dF_q}\left(\chi\circ N_{k/\dF_q},\psi\right) = \varepsilon(k) \tau(\chi,\psi)^N \text{.} 
\end{equation*}

If $k$ is a field, $\varphi$ is a circular permutation of $I$,
$\varepsilon(k) = (-1)^{N+1}$, and we recover the Hasse-Davenport identity: 
\[
  \tau\left(\chi\circ N_{k/\dF_q},\psi\circ \tr_{k/\dF_q}\right) = (-1)^{N+1} \tau_{\dF_q}\left(\chi\circ N_{k/\dF_q},\psi\right) = \tau(\chi,\psi)^N \text{.} 
\]





\begin{lemma_}\label{VI:4-13}
For $\chi\circ \tau=(\chi_i)_{i\in I}$ as in \ref{VI:4-9}(i), the
following conditions are equivalent 
\begin{enumerate}[\indent (i)]
  \item $\chi| \dF_q^\times$ is trivial
  \item The product of the $\chi_i$ is trivial
  \item $\sF(\chi) = \sK_n((\chi_i))$ is the inverse image of a (unique)
    sheaf on $V^\ast/\dG_m$.
  \item The inverse image of $\sF(\chi)$ on $\dG_m$ (mapped into
    $V^\ast$ via the structural morphism $\dF_q \to k$) is trivial.
\end{enumerate}
\end{lemma_}

(i)$\Rightarrow$(ii). We regard $\chi$ as a character of
$V^\ast/\dG_m(\dF_q)=k^\times/\dF_q^\times$, and we take $\sF(\chi)$ on
$V^\ast/\dG_m$.

(ii)$\Rightarrow$(iii). Likewise, we regard $(\chi_i)$ as a character
of $(V^\ast/\dG_m)_n$. Uniqueness in (iii) follows from $V^\ast$ being
a bundle with connected fibres over $V^\ast/\dG_m$, and (iii)$\Rightarrow$(iv)
is trivial.

(iv)$\Rightarrow$(i), (ii). This inverse image is
$\sF(\chi|\dF_q^\times)$ and $\sK_n(\prod \chi_i)$. 





\subsection{}\label{VI:4-14}

Let $k$ be an étale algebra over $\dF_q$, of
dimension $N+1$ and $\chi$ a non-trivial character of $k^\times$, trivial on
$\dF_q^\times$. The Jacobi sum $J(\chi)$ is defined by 
\begin{equation*}\tag{4.14.1}\label{VI:eq:4-14-1}
  J(\chi) = (-1)^{N-1} \sum_{\substack{x\in k^\times/\dF_q^\times \\ \tr(x)=0}} \chi^{-1}(x) \text{.} 
\end{equation*}

We have between Gauss sums and Jacobi sums the following identity. 





\begin{proposition_}\label{VI:4-15}
For $\chi$ as above and $\psi$ a non-trivial additive character of
$\dF_q$, we have 
\[
  q J(\chi) = \tau_{\dF_q}(\chi,\psi) \text{.} 
\]
\end{proposition_}

In the special case where $k=\dF_q^n$, this formula is rewritten by
\eqref{VI:eq:4-5-3} 
\begin{equation*}\tag{4.15.1}\label{VI:eq:4-15-1}
  q J(\chi) = \prod_i \tau(\chi_i,\psi) \text{,} 
\end{equation*}
more written in the form 
\begin{align*}
  \chi_0(-1) J(\chi) &= (-1)^{N-1} \sum_{\substack{x_1,\dots,x_N\in \dF_q^\times \\ \sum x_i = 1}} \prod_{i=1}^N \chi_i^{-1}(x_i) = \tau(\chi_0^{-1},\psi)^{-1} \prod_{i=1}^N \tau(\chi_i,\psi) \\
  \chi_0^{-1} &= \prod_{i=1}^N \chi_i 
\end{align*}





\begin{proof}[Proof of \ref{VI:4-15}]
\[
  \tau_{\dF_q}(\chi,\psi) = (-1)^{N+1} \sum_{x\in k^\times} \chi(x)^{-1} \psi \tr(x) = (-1)^{N+1} \sum_{x\in k^\times/\dF_q^\times} \chi(x)^{-1} \sum_{\lambda\in \dF_q^\times} \psi(\lambda \tr x) \text{.}
\]
The sum $\sum_{\lambda\in \dF_q^\times} \psi(\lambda \tr x)$ is $q-1$ if
$\tr(x)=0$, and $-1$ if $\tr(x)\ne 0$. Hence, 
\[
  \tau_{\dF_q}(\chi,\psi) = (-1)^{N+1} \left(q\sum_{x\in k^\times/\dF_q^\times} \chi(x)^{-1} - \sum_{x\in k^\times/\dF_q^\times} \chi(x)^{-1}\right) \text{.}
\]
The second term on the right-hand side is zero, as the sum of the values of a
character, and \ref{VI:4-15} follows. 
\end{proof}

Proposition \ref{VI:4-15} transposes as follows into cohomology: 





\begin{proposition_}\label{VI:4-16}
The cohomology of $W^\ast/\dG_m$ (\ref{VI:4-6}, \ref{VI:4-8}) with coefficients
in $\sF(\chi^{-1})$ satisfies
\begin{enumerate}[\indent (i)]
  \item $\h_c^i=0$ for $i\ne N-1$, and $\dim \h_c^{N-1}{}=1$.
  \item On $\h_c^{N-1}$, $F^\ast$ is multiplication by $J(\chi)$.
  \item $\h_c^{N-1}$ is canonically isomorphic to
    $\h_c^{N+1}\left(V^\ast,\sF(\chi^{-1},\psi\circ \tr_{k/\dF_q})\right)(1)$.
\end{enumerate}
\end{proposition_}

Assertion (ii) follows from (i) and the trace formula; (i) and (iii)
follow from \ref{VI:4-11} and the following more geometric statement. 





\begin{proposition_}\label{VI:4-17}
Let $(\chi_i)_{i\in I}$ be a family of everywhere non-trivial characters of
$\dmu_n(\dF)$, of product $1$. Then, 
$\h_c^{i-1}\left(W^\ast/\dG_m,\sK_m((\chi_i)_{i\in I})\right)$ is
canonically isomorphic to 
\newline
$\h_c^{i+1}\left(V^\ast,\sK_m((\chi_i)_{i\in I})\otimes \sF(\psi\tr_{k/\dF_q})\right)(1)$. 
\end{proposition_}

The first line of \ref{VI:4-15} becomes ($\pi$ as in
\eqref{VI:eq:4-6-1}). 
\begin{equation*}\tag{4.17.1}\label{VI:eq:4-17-1}
  \eR^i \pi_!\left(\sK_n((\chi_i)_{i\in I})\otimes \sF(\psi \tr_{k/\dF_q})\right) = \sK_n((\chi_i)_{i\in I}) \otimes \eR^i \pi_! \sF(\psi\tr_{k/\dF_q}) 
\end{equation*}
(on $V^\ast/\dG_m$). Let us compute the sheaves
$\eR^i\pi_! \sF(\psi\tr_{k/\dF_q}) = \eR^i \pi_! T^\ast \sF(\psi)$. Let
$v_0:\widetilde V_0 \to V_0$ be the blow-up of $V_0$ at $\{0\}$. In the
diagram 
\begin{equation*}\tag{4.17.2}\label{VI:eq:4-17-2}
\xymatrix{
  V_0\setminus \{0\} \ar@{^{(}->}[r] \ar[dr]_-\pi 
    & \widetilde V_0 \ar[r]^-{v_0} \ar[d]^-{\bar\pi} 
    & V_0 \\
  & P_0 
}
\end{equation*}
$\pi$ is a fibration with punctured-line fibre, $\widetilde V_0$
is the corresponding line bundle, and $Z_0=v_0^{-1}(0)$ its $0$-section.
The sheaf $T^\ast\sF(\psi)$ on $V-0\setminus \{0\}$ extends to
$(T v_0)^\ast \sF(\psi)$ on $\widetilde V_0$, and the restriction of this sheaf
to $Z_0$ is the constant sheaf $E_\lambda$, for $T_0 v_0$ vanishes on
$Z_0$. The long exact proper-support cohomology sequence therefore
reads 
\begin{equation*}\tag{4.17.3}\label{VI:eq:4-17-3}
\xymatrix{
  \ar[r] 
    & \eR^i \pi_! T^\ast \sF(\psi) \ar[r] 
    & \eR^i \bar\pi_! (T v)^\ast \sF(\psi) \ar[r] 
    & (E_\lambda\text{ for $i=0$, $0$ otherwise}) \ar[r] 
    & 
}
\end{equation*}

The inverse image $\bar\pi^{-1}(Q_0)$ is the blow-up $\widetilde W_0$ of
$W_0$ at $0$; $T_0 V_0$ vanishes on $\widetilde W_0$; we therefore have
$(T v)^\ast \sF(\psi)=E_\lambda$ on $\pi^{-1}(Q_0)$ and, $\bar\pi$ being a
line bundle 
\begin{equation*}\tag{4.17.4}\label{VI:eq:4-17-4}
  \eR^i \bar\pi_! (T v)^\ast \sF(\psi)|Q = 
    \begin{cases}
      0 & \text{for $i\ne 2$} \\
      E_\lambda(-1) & \text{for $i=2$.} 
    \end{cases} 
\end{equation*}

If $x\in P$, $x\notin Q$, the line $D=\pi^{-1}(x)$ is sent
isomorphically onto $\dG_a$ by $T v$; we therefore have
\ref{VI:2-7}*) 
\begin{align} \notag 
  \h_c^\bullet(D,(T v)^\ast \sF(\psi)) &= \h^\bullet(\dG_a,\sF(\psi)) = 0 && \text{and} \\
  \tag{4.17.5}\label{VI:eq:4-17-5}
  \eR^i \bar\pi_! (T v)^\ast \sF(\psi) 
    &= \begin{cases} 
         0 & \text{for $i\ne 2$} \\
         E_\lambda(-1)_Q & \text{for $i=2$} 
       \end{cases} 
\end{align}

Combining \eqref{VI:eq:4-17-3} and \eqref{VI:eq:4-17-5}, we finally find 
\begin{equation*}\tag{4.17.6}\label{VI:eq:4-17-6}
  \eR^i \pi_! T^\ast\sF(\psi) = 
    \begin{cases}
      0 & \text{for $i\ne 1,2$} \\
      E_\lambda & \text{for $i=1$} \\
      E_\lambda(-1)_Q & \text{for $i=2$.} 
    \end{cases}
\end{equation*}





\subsection{}\label{VI:4-18}

Let us compute the proper-support cohomology of the sheaf
$\sK_n((\chi_i))\otimes T^\ast \sF(\psi)$ on $V^\ast$ using the
Leray spectral sequence of $\pi:V^\ast\to V^\ast/\dG_m$. Applying
\eqref{VI:eq:4-17-1} and \eqref{VI:eq:4-17-6}, we find as initial terms
the $E_2^{p,1} = \h_c^p(V^\ast/\dG_m,\sK_n(\chi_i)) = 0$ (for $(\chi_i)\ne 0$)
and the $E_2^{p,2} = \h_c^p(W^\ast/\dG_m,\sK(\chi_i))(-1)$.

The spectral sequence reduces to an isomorphism, and \ref{VI:4-17}
follows. 





\subsection{}\label{VI:4-19}

The sheaves $\sK_n((\chi_i)_{i\in I})$ make sense over any field,
indeed over any base scheme. This will allow us to generalise
\ref{VI:4-16}(i). With the notations of \ref{VI:4-6}, suppose $M$ finite
étale everywhere of rank $N$ over $A$, and let $n$ be an integer invertible in $A$.
We denote by $a$ the projection from $W^\ast/\dG_m$ to $\spec(A)$. Let also $\chi$
be a character (morphism of sheaves)
$\chi:(V^\ast/\dG_m)_n \to E_\lambda^\times$, and $\sK_n(\chi)$ the corresponding
sheaf. Locally for the étale topology, we may regard $\chi$
as a family of characters $(\chi_i)_{i\in I}$ of $\dmu_n$, of trivial product.





\begin{proposition_}\label{VI:4-20}
\begin{enumerate}[(i)]
  \item If $\chi$ is (at every point) non-trivial, we have
    $\eR^i a_! (\sK_n(\chi)) = 0$ for $i\ne N-1$, and
    $\eR^{N-1} a_!(\sK_n(\chi))$ is smooth, of rank $1$.
  \item An automorphism $\sigma$ of $M$ which respects $\chi$ acts on this
    $\eR^{N-1} a_!$ by multiplication by the signature $\varepsilon(\sigma)$
    of $\sigma$, viewed as a permutation of $I$.
  \item If the $\chi_i$ are all non-trivial, we have $\eR a_! \iso \eR a_\ast$.
\end{enumerate}
\end{proposition_}
\begin{proof}[Proof]
The scheme $W^\ast/\dG_m$ is the complement of a relative normal-crossings
divisor in the scheme $Q$, proper and smooth over $\spec(A)$, and
$\sK_n(\chi)$ is locally constant on $W^\ast/\dG_m$, tamely ramified
at infinity. It follows that the $\eR^i a_!$ and
$\eR^i a_\ast$ are smooth, of formation compatible with any base
change. A standard argument then reduces us to supposing that $A$ is a
finite field, and (i) follows from \ref{VI:4-17} and \ref{VI:4-11}(i). To
prove (ii), we use that the action of $\sigma$ is compatible with
isomorphism \ref{VI:4-17}, and \ref{VI:4-12}. For (iii), we note that if the
$\chi_i$ are all non-trivial, then $\sK_n(\chi)$ on
$W^\ast/\dG_m\subset Q$ is ramified along each divisor at infinity,
and \ref{VI:1-19-1}.
\end{proof}










\section{Hecke characters}\label{VI:5}





\subsection{}\label{VI:5-1}

Let $F$ be a number field (of finite degree over $\dQ$) and $k$ a field of
characteristic $0$. Here are various equivalent ways of saying what an
\emph{algebraic} homomorphism $\alpha:F^\times \to k^\times$ is.
\begin{enumerate}[(i)]
  \item Let $\{e_i\}$ be a basis of $F$ over $\dQ$. The homomorphism $\alpha$ is
    algebraic if it is given by a formula
    $\alpha(\sum x^i e_i) = A(x^i)$, $A\in k(X^i)$.
\end{enumerate}

This means that $\alpha$ coincides on $F^\times$ with a rational
map defined over $k$ from $\res_{F/\dQ}(\dG_m)$ to $\dG_m$. By
Zariski density of $F^\times$, and the fact that a birational homomorphism
is everywhere defined, such a map is a homomorphism of group
schemes: 
\begin{enumerate}[(i)]
\setcounter{enumi}{1}
  \item $\alpha$ is induced by a homomorphism of $k$-group schemes
    \[
      \res_{F/\dQ}(\dG_m)\otimes_\dQ k \to \dG_m \text{.}
    \]
\end{enumerate}

If $k$ is a number field, the adjunction property of Weil restriction
of scalars $R$ shows that this is equivalent to 
\begin{enumerate}[(i)]
\setcounter{enumi}{2}
  \item ($k$ a number field) $\alpha$ is induced by a homomorphism of
    $\dQ$-group schemes: $\res_{F/\dQ}(\dG_m) \to \res_{k/\dQ}(\dG_m)$.
\end{enumerate}

Let $\bar k$ be an algebraic closure of $k$, and $I=\hom(F,\bar k)$. Over
$\bar k$, the character group of $\res_{F/\dQ}(\dG_m)$ is $\dZ^I$, with
basis the embeddings of $F$ into $\bar k$. The characters defined over $k$
are those invariant under $\gal(\bar k/k)$; they can be described either as
those sending $F^\times$ into $k^\times\subset \bar k$, or in terms of the
orbits of $\gal(\bar k/k)$ in $I$, themselves corresponding to the factors of
$F\otimes k$. 
\begin{enumerate}[(i)]
\setcounter{enumi}{3}
  \item $\alpha$ is of the form
    $\alpha=\prod_{\omega\in I} \omega^{n_\omega}$. The permitted families of exponents
    $(n_\omega)$ are those with $n_\omega=n_{\omega'}$ for
    $\omega$ and $\omega'$ in the same orbit of $\gal(\bar k/k)$. They are
    also those with $\prod_\omega (x)^{n_\omega}\in k$ for every
    $x\in F$.
  \item Put $F\otimes k=\prod_{j\in J} F_j$, the $F_j$ being fields.
    Then $\alpha$ is written $\alpha=\prod N_{F_j/k}^{m_j}$.
\end{enumerate}

In the special case where $k$ contains a normal closure of $F$, these
expressions become: $\alpha=\prod \omega^{n_\omega}$, where $\omega$
runs over the $[F:\dQ]$ embeddings of $F$ into $k$. 





\subsection{}\label{VI:5-2}

Suppose that $k$ is a number field, and let $\alpha$ be an algebraic
homomorphism from $F^\times$ to $k^\times$. There then exists one and only one
homomorphism, still denoted $\alpha$, from the group $I(F)$ of fractional
ideals of $F$ to that of $k$, such that $\alpha((x))=(\alpha(x))$.
Uniqueness follows from every ideal having a power which is a
principal ideal, and from $I(k)$ being torsion-free. To prove
existence, we use for example \ref{VI:5-1}(v): we have
$\alpha(a) = \prod N_{F_j/k}^{m_j}((a))$. 





\subsection{}\label{VI:5-3}

Recall the definition of algebraic Hecke characters, called
Weil Hecke characters (or: gr\"ossencharaktere) of type $A_0$. Let $F$
be a number field, $\fm$ an ideal of $F$ (i.e., of the ring of integers of
$F$), $I_\fm$ the group of fractional ideals of $F$ prime to $\fm$,
and $k$ a field of characteristic $0$. A homomorphism
$\chi:I_\fm \to k^\times$ is an \emph{algebraic Hecke character} (of
conductor $\leqslant \fm$) if there exists an algebraic homomorphism
$\chi_\text{alg}:F^\times \to k^\times$ satisfying 
\begin{equation*}\tag{$*$}\label{VI:eq:5-3-1}
\text{For $x\in F^\times$, prime to $\fm$, totally positive and
$\equiv 1\pmod\fm$, we have $\chi((x))=\chi_\text{alg}(x)$.}
\end{equation*}

By density of the set of $x$ in \eqref{VI:eq:5-3-1}, $\chi_\text{alg}$
is entirely determined by $\chi$. It is the \emph{algebraic part}
of $\chi$. If $\chi((x))=\chi_\text{alg}(x)$ for $x$ totally positive and
$\equiv 1\pmod{\fm'}$, $\chi$ extends to a Hecke character of
conductor $\leqslant \inf(\fm,\fm')$: $I_{\fm+\fm'}\to k^\times$. We
will identify Hecke characters which coincide on their common domain of
definition, and we call \emph{conductor} of $\chi$ the smallest $\fm'$
such that $\chi$ is of conductor $\leqslant \fm'$. 





\subsection{Remark}\label{VI:5-4}

If an algebraic Hecke character $\chi$ takes its values in a
subfield $k'$ of $k$, it is already a Hecke character with values in
$k'$: it suffices to see that $\chi_\text{alg}:\res_{F/\dQ}(\dG_m) \to \dG_m$
is already defined over $k'$, and this follows from Zariski density
in $\res_{F/\dQ}(\dG_m)$ of the set of totally positive $x$
$\equiv 1\pmod\fm$. One can always take for $k'$ a subfield of $k$ of
finite degree over $\dQ$. 





\subsection{}\label{VI:5-5}

If $\varepsilon$ is a totally positive unit $\equiv 1\pmod\fm$, we have
$\chi_\text{alg}(\varepsilon) = \chi((\varepsilon)) = 1$. The homomorphism
$\chi_\text{alg}$ therefore factors through the quotient $T_\fm$ of
$\res_{F/\dQ}(\dG_m)$ by the Zariski closure of the group
$E_\fm\subset F^\times$ of totally positive units $\equiv 1\pmod\fm$.

By Serre \cite[II.3]{se68}, if $k$ is an algebraic closure of
$\dQ$, and $\fm$ is large enough, the characters $\prod \omega^{n_\omega}$
of $\res_{F/\dQ}(\dG_m)$ of the form $\chi_\text{alg}$ are characterised
as follows: there must exist an integer $N$, the \emph{weight} of $\chi$ (or of
$\chi_\text{alg}$), such that for every element $\sigma$ of $\gal(k/\dQ)$
conjugate to complex conjugation, we have $n_\omega+n_{\sigma \omega}=N$. If
$F_1$ is the largest subfield of $F$ which is a totally imaginary quadratic
extension of a totally real field $F_1'$, this amounts to
saying that $\chi_\text{alg}$ is of the form $\chi_1\circ N_{F/F_1}$, and that
$\chi_1|F_1=(N_{F_1/\dQ})^N$.

If $\chi$ is of weight $N$, for every ideal $\fa$ of $F$, $\chi(\fa)$ is an
algebraic number all of whose conjugates have absolute value
$N(\fa)^{N/2}$ \cite[II.3 prop.2]{se68}. 





\subsection{}\label{VI:5-6}

For $k$ a fixed number field, additional conditions are
imposed on $\chi_\text{alg}$. For every ideal $\fa$ of $F$ prime to the
conductor, we indeed have 
\begin{equation*}\tag{5.6.1}\label{VI:eq:5-6-1}
  \chi_\text{alg}(\fa) = (\chi(\fa)) \text{,} 
\end{equation*}
so that $\chi_\text{alg}(\fa)$ is principal. Since the group of
fractional ideals of $k$ is torsion-free, it suffices to prove the
$n$-th power of \eqref{VI:eq:5-6-1}, $n\ne 0$ suitable; this allows us
to replace $\fa$ by $\fa^n$, hence to suppose $\fa=(x)$, with $x$ totally
positive $\equiv 1\pmod\fm$. It only remains to use the definitions.

This, together with \ref{VI:5-5}, shows that $\chi_\text{alg}$ determines the norm
of the $\chi(\fa)$ at all places of $k$. 





\subsection{}\label{VI:5-7}

Serre's group $S_\fm$ could be characterised as being the
group of multiplicative type whose character group (over any
field) is the group of algebraic Hecke characters of conductor
$\leqslant \fm$. (cf. \cite[II 2.1 et 2.2]{se68}). Its relation with
$\ell$-adic representations is explained in \cite[II 2.3]{se68}. 





\begin{theorem_}[{\cite{se68}}]\label{VI:5-8}
Let $\chi$ be an algebraic Hecke character of $F$ with values in $E_\lambda$, for
$E_\lambda$ a finite extension of $\dQ_\ell$. There then exists one (and only one)
homomorphism $\chi_\lambda:\gal(\bar F/F)^\textnormal{ab} \to E_\lambda$, such
that
\begin{enumerate}[\indent (i)]
  \item $\chi_\lambda$ is unramified outside the conductor $\ff$ of
    $\chi$ and $\ell$.
  \item For $\fp$ a prime ideal of $F$ prime to $\ff$ and to $\ell$,
    and $F_\fp\in \gal(\bar F/F)^\textnormal{ab}$ the geometric Frobenius
    at $\fp$, we have 
    \[
      \chi_\lambda(F_\fp) = \chi(\fp) \text{.}
    \] 
\end{enumerate}
\end{theorem_}





\subsection{}\label{VI:5-9}

In the rest of this section, we will give a criterion for a
$\lambda$-adic representation to come in this way from a Hecke character,
and, in the next section, we will apply it to Jacobi sums. We
will thus recover, somewhat generalised, results of Weil
\cite{we52,we74-2}, with the second part of Weil's proof
replaced by an $\ell$-adic cohomology argument. 





\begin{theorem_}\label{VI:5-10}
Let $F$ and $k$ be two number fields, $\lambda$ a place of $k$ of
characteristic $\ell$, and
$\chi_\lambda:\gal(\bar F/F)^\textnormal{ab}\to k_\lambda^\times$ a
homomorphism, unramified outside $\ell$ and a finite set $S$ of
places. Suppose there exist a set $T$ of places (containing $S$ and the
places above $\ell$), of density $0$, and an algebraic homomorphism
$\chi_0:F^\times \to k^\times$ such that
\begin{enumerate}[(i)]
  \item For $\bar k$ an algebraic closure of $k$,
    $\chi_0:F^\times \to k^\times \hookrightarrow \bar k^\times$ is of the type
    considered in \ref{VI:5-5}, of weight $N$;
  \item For $\fp$ prime $\notin T$, $\chi_\lambda(F_\fp)$ lies in
    $k^\times$, $(\chi_\lambda(F_\fp))=\chi_0(\fp)$, and all the complex
    conjugates of $\chi_\lambda(F_\fp)$ have absolute value $(N \fp)^{N/2}$.
\end{enumerate}

Then $\chi_\lambda$ is defined (\ref{VI:5-8}) by an algebraic Hecke
character of $F$ with values in $k$, of algebraic part $\chi_0$.
\end{theorem_}

Let $\chi'$ be an algebraic Hecke character with values in a
finite Galois extension $k'$ of $k$, of algebraic part $\chi_0$
(\ref{VI:5-5}). Enlarging $T$ if necessary, we may suppose that the conductor of
$\chi'$ is supported in $T$. Let $\lambda'$ be a place of $k'$ above
$\lambda$, $\chi_{\lambda'}$ the composite
$\gal(\bar F/F)^\text{ab}\to k_\lambda^\times \to {k'}_{\lambda'}^\times$, and
$\chi_{\lambda'}':\gal(\bar F/F)^\text{ab}\to {k'}_{\lambda'}^\times$ defined
by $\chi'$ (\ref{VI:5-8}). Put
$\varepsilon = \chi_{\lambda'} \cdot {\chi'}_{\lambda'}^{-1}$.

Hypothesis (ii), and \ref{VI:5-5}, \ref{VI:5-6}, \ref{VI:5-8} ensure that
for $\fp\notin T$, we have $\varepsilon(F_\fp)\in {k'}^\times$, and that, in all
completions of $k'$, this number has norm $1$. It follows that
$\varepsilon(F_\fp)$ belongs to the (finite) group $\dmu$ of roots of unity
of $k'$. By the Cebotarev density theorem, and the density
hypothesis on $T$, the $F_\fp$ ($\fp\notin T$) are dense in
$\gal(\bar F/F)^\text{ab}$, so that $\varepsilon(\sigma)\in \mu$ for every
$\sigma\in \gal(\bar F/F)^\text{ab}$: the character $\varepsilon$ is a
character of finite order $\gal(\bar F/F)^\text{ab}\to \mu$. By
class field theory, it corresponds to a character of finite order
of the idèle class group of $F$, i.e. to an algebraic Hecke
character of trivial algebraic part. Correcting $\chi'$ by this
character, we reduce to supposing that
$\chi_{\lambda'}=\chi_{\lambda'}'$, and it remains to show that $\chi'$, a
priori with values in $k'$, in fact takes values in $k$. If
$\sigma\in \gal(k'/k)$, $\chi'$ and ${\chi'}^\sigma$ coincide on every prime
ideal $\fp\notin T$, since
$\chi'(\fp) = \chi'_{\lambda'}(F_\fp) = \chi_{\lambda'}(F_\fp) \in k^\times$.
The characters $\chi_{\lambda'}'$, and ${\chi_{\lambda'}'}^\sigma$, therefore
coincide on a dense subset of $\gal(\bar F/F)^\text{ab}$, hence are equal, so
that $\chi'={\chi'}^\sigma$: $\chi'$ takes values in $k$. 










\section{Hecke characters defined by Jacobi sums}\label{VI:6}

The idea of this section is as follows. Let a Gauss sum
$\tau_{\dF_q}(\chi,\psi)$ be given. It, and the corresponding sheaf
$\sF(\chi^{-1}\psi)$ are doubly linked to the characteristic $p$, and to the
base field $\dF_q$:
\begin{enumerate}[\indent a)]
  \item for there appears the additive character $\psi$, and the sheaf
    $\sf(\psi)$;
  \item for $\sF(\chi)$ is defined from the Lang exact sequence 
    \[\xymatrix{
      0 \ar[r] 
        & k^\times \ar[r] 
        & V^\ast \ar[r] 
        & V^\ast \ar[r] 
        & 0 \text{.} 
    }\]
\end{enumerate}
Two difficulties if one wants to lift to a number field $F$ the
corresponding cohomological constructions, and to find
$\ell$-adic representations of $\gal(\bar F/F)$ where the Frobenius
eigenvalues are Gauss sums.

If $\chi$ is trivial on $\dF_q^\times$, the sum $\tau_{\dF_q}(\chi,\psi)$
is independent of $\psi$, and $q$ times a Jacobi sum. The latter is
linked to the cohomology of $W^\ast/\dG_m$, with values in a sheaf
deduced from $\chi$: $\psi$, and difficulty a), have disappeared. To
resolve b), it suffices to describe the sheaves used by means of
K\"ummer exact sequences (which make sense in any characteristic). This is
possible thanks to \ref{VI:4-9}.

Once the $\ell$-adic representations are found, we use
\ref{VI:5-10} and Stickelberger's theorem to show that they
come from algebraic Hecke characters. 





\subsection{}\label{VI:6-1}

If $E$ is a field, of algebraic closure $\bar E$, and $\lambda$ is
a character of order dividing $n$ of $\mu_d(\bar E)$, we denote by $_n\lambda$ the
character of $\mu_n(\bar E)$ such that $\lambda$ and $_n \lambda$ induce the
same character of $\widehat\dZ(1)_{\bar E} = \varinjlim \mu_n(\bar E)$. Likewise
for $\lambda$ a character of $\widehat\dZ(1)_{\bar E}$.

Let $F$ and $k$ be two number fields, $\bar F$ an algebraic closure of
$F$, $I$ a finite set equipped with an action of $\gal(\bar F/F)$ and
$\lambda=(\lambda_i)_{i\in I}$ a family of characters
$\lambda_i:\widehat\dZ(1)_{\bar F} \to k^\times$. We suppose that
$\lambda_i\ne 1$ ($i\in I$), that the product of the $\lambda_i$ is trivial and
that the family $\lambda$ is defined over $F$, i.e. that for every
$\sigma\in \gal(\bar F/F)$, we have
$\lambda_{\sigma(i)}=\lambda_i\circ \sigma^{-1}$.

The Galois set $I$ is the set of homomorphisms into $\bar F$ of an
étale algebra $E$ over $F$. If $C$ is the set of orbits of
$\gal(\bar F/F)$ in $I$, we have a decomposition of $E$ as a product of fields
$E=\prod_{\alpha\in C} E_\alpha$, where $\alpha$ identifies with the set of
embeddings of $E_\alpha$ into $\bar F$. If the characters $\lambda_i$, for
$i\in \alpha$, are of order $d_\alpha$, the field $E_\alpha$ contains the
$d_\alpha$-th roots of $1$ and there exists
$\lambda_\alpha:\mu_{d_\alpha}(E_\alpha) \hookrightarrow k^\times$ such that, if
$i\in \alpha$ corresponds to the embedding $\sigma_i$ of $E_\alpha$ into $\bar F$,
we have $_{d_\alpha} \lambda_i = \lambda_\alpha \sigma_i^{-1}$. 





\subsection{}\label{VI:6-2}

Let $n>0$ be a multiple of the $d_\alpha$, $\Sigma$ the set of places of $F$
where $E/F$ ramifies or which divide $n$, $\mu$ a place of $k$ of
residue characteristic $\ell$, and $A$ the ring of elements of $F$
integral outside $\Sigma$ and $\ell$. The integral closure $M$ of $A$
is finite étale of rank $N=|I|$ over $A$. With the notations of \ref{VI:4-6},
\ref{VI:4-8}, the $_n \lambda_i$ define 
\[
  _n\lambda:(V^\ast/\dG_m)_n \to k^\times \subset k_\mu^\times \text{.} 
\]

Let $a$ be the projection from $W^\ast/\dG_m$ to $\spec(A)$. By
\ref{VI:4-20}, the sheaf $\eR^{N-2} a_! \sK_n(_n \lambda)$ is smooth of rank
one; it defines an $\ell$-adic representation
\[
  j_\mu[\lambda]:\gal(\bar F/F)^\text{ab} \to k_\mu^\times
\]
(or simply $j_\mu$) unramified outside $\Sigma$ and $\ell$. 





\subsection{}\label{VI:6-3}

Let us compute $j_\mu(F_\fp)$. By \ref{VI:4-20}, it suffices to do so
after reduction modulo $\fp$. Let therefore $f=A/\fp$, $e=M/\fp M$, and
$\bar f$ the algebraic closure of $f$ defined by a place of $\bar F$
above $\fp$. We still have $I=\hom_f(e,\bar f)$. After reduction,
$_n\lambda$ still admits a description \ref{VI:6-1}:
\begin{enumerate}[a)]
  \item Let $D$ be the set of orbits of $\gal(\bar f/f)$ in $I$. The
    decomposition of $e$ as a product of fields is written
    $e=\prod_{\beta\in D} e_\beta$, where the $\sigma_i$ ($i\in \beta$)
    identify with the embeddings of $e_\beta$ into $F$. For $\beta\in D$, denote
    by $\alpha(\beta)$ the element of $C$ containing $\beta$ and
    $d_\beta=d_{\alpha(\beta)}$. Let 
    \[\xymatrix{
      \lambda_\beta : \mu_{d_\beta}(e_\beta) 
        & \ar[l]_-\sim \mu_{d_\beta}(E_{\alpha(\beta)}) \ar[r]^-{\lambda_{d(\beta)}} 
        & k^\times \text{.} 
    }\]
  \item If $i\in I$ corresponds to the embedding $\sigma_i$ of $e_\beta$ into
    $\bar f$, the reduction
    $_{d\beta} \bar\lambda_i:\mu_{d\beta}(\bar f) \to k^\times$ of
    $_{d\beta} \lambda_i$ is
    $\lambda_\beta \circ \sigma_i^{-1}$.
\end{enumerate}

Let $q_\beta$ be the number of elements of $e_\beta$,
$\chi_\beta=_{(q_{\beta^{-1}})} \lambda_\beta$, and
$\chi:e^\times \to k^\times$ of coordinates the $\chi_\beta$. On the
affine space over $A/\fp$ defined by $e$ (\ref{VI:4-6}), or rather on the torus
$e_\beta^\times$, we have isomorphisms 
$\sK_n((_n\bar\lambda_i)_{i\in\beta}) = \sK_{d_\beta}((\lambda_\beta \circ \sigma_i^{-1})_{i\in\beta}) =\sK_{q_{\beta^{-1}}}((\chi_\beta\circ \sigma_i^{-1})_{i\in I}) = \sF(\chi_\beta)$ (\ref{VI:4-9}). 
On $e^\times$, we therefore have $\sK_n((_n\bar\lambda_i)_{i\in \beta}) = \sF(\chi)$
and this isomorphism descends to the reduction mod $\fp$ of $V^\ast/\dG_m$.
Hence





\begin{proposition_}\label{VI:6-4}
With the notations of \ref{VI:6-2} and \ref{VI:6-3}, $j_\mu(F_\fp)$ is the
Jacobi sum $J(\chi)$.
\end{proposition_}

Let us rewrite the formulas linking $(\lambda_i)_{i\in I}$ to
$(\chi_\beta)_{\beta\in B}$: for $i\in \beta\subset \alpha$, 
\begin{align*}
  _{d_\alpha} \lambda_i &= \lambda_\alpha\circ \sigma_i^{-1} && \lambda_\alpha:\mu_{d_\alpha}(E_\alpha) \to k^\times \\
  \lambda_\beta &= \text{reduction of $\lambda_\alpha$} && :\mu_{d_\alpha}(e_\beta) \to k^\times \\
  \chi_\beta &= _{q_{\beta^{-1}}}\lambda_\beta: x\mapsto \lambda_\beta(x^{(q_\beta-1)/d_\alpha}) && :e_\beta^\times \to k^\times 
\end{align*}





\begin{theorem_}\label{VI:6-5}
The representation $j_\mu$ is defined by an algebraic Hecke
character of $F$ with values in $k^\times$.
\end{theorem_}

We will check this using \ref{VI:5-10}. We here rejoin the
proof of Weil \cite{we52,we74-2}, through the use of Stickelberger's
theorem, and I will give only a few indications. 





\subsection{}\label{VI:6-6}

Given $j:\gal(\bar F/F) \to k_\lambda^\times$, and a finite extension
$F'/F$, one checks using \ref{VI:5-10} that $j$ is defined by a
Hecke character if and only if $j|\gal(\bar F/F')$ is. This
remark allows us in \ref{VI:6-5} to reduce to the case where Galois acts
trivially on $I$, and where $F$ contains the $n$-th roots of
unity. One can then reduce to the case where $F=\dQ(\sqrt[n]{1})$, then
take $k=F$. In this case, each $_n \lambda_i$ is defined by
$a_i\in \dZ/n$: it is $x\mapsto x^{a_i}:\mu_n \to k^\times$. We have
$\sum a_i=0$, and $a_i\ne 0$. 





\subsection{}\label{VI:6-7}

In a number field, the places of degree $1$ always form a
set of density $0$. Applying \ref{VI:5-10}, we can therefore neglect
the others, and use Stickelberger only for a prime field. Let therefore
$p$ be prime, $0<a<p-1$, $\widetilde x:\dF_p^\times \to \dQ_p^\times$ the
``multiplicative lift'' character and
$g=-\sum_{x\in \dF_p^\times} \widetilde x^{-a} \zeta^z\in \dQ_p(\zeta)$, for
$\zeta$ a $p$-th root of $1$. Let us still denote by $x$ the integer in $[0,p-1]$
lifting $x$, and expand $\zeta^x=(1+\pi)x$ by the binomial formula.
We find $g=-\sum_{0=0}^{p-1} A_i\pi^i$ with $A_i\in \dZ_p$, of reduction
$\mod p$ the sum $\sum_{x\in \dF_p^\times} x^{-a} \binom x i$. For $b$ not
divisible by $p-1$, we have $\sum_{x\in \dF_p^\times} x^b=0$. If $i<a$,
$\binom x i$ is a polynomial in $x$ of degree $i<a$, whence
$A_i\equiv 0\pmod p$. Likewise, $A_a\equiv \frac{p-1}{a!} \pmod p$, and
$g\sim \frac{\pi^a}{a!}$: the valuation of $g$ is given by
$v(g)/v(p) = \frac{a}{p-1}$. 





\subsection{}\label{VI:6-8}

Let $\sigma_i$ ($i\in (\dZ/n)^\times$) be the element of
$\gal(\dQ(\sqrt[n]{1})/\dQ)$ with $\sigma_i(\zeta)=\zeta^i$ for $\zeta^n=1$,
and let us denote additively the group of algebraic homomorphisms
$\dQ(\sqrt[n]{1})^\times \to \dQ(\sqrt[n]{1})^\times$. For $\lambda$ as in
\ref{VI:6-6}, defined by a family $(a_i)_{i\in I}$ of integers $\mod n$, non-
zero and of zero sum \ref{VI:5-10}, \eqref{VI:eq:4-15-1} and \ref{VI:6-7}
show that $j_\mu[\lambda]$ is defined by an algebraic Hecke
character $j[a]$ and that, if $N$ denotes the norm, its algebraic part is
given by 
\begin{equation*}\tag{6.8.1}\label{VI:eq:6-8-1}
  (N j[a])_\text{alg} = \sum_{i\in (\dZ/n)^\times} \sum_{j\in I} \left \lfloor \frac{i a_j}{n} \right\rfloor \sigma_i^{-1} 
\end{equation*}
where $\lfloor \cdot \rfloor$ is the integral part. 





\subsection{}\label{VI:6-9}

The most interesting case is that where $F\subset \dQ(\sqrt[n]{1})$, and
where the family of $_n \lambda_i$ is constructed as follows:
\begin{enumerate}[a)]
  \item We take a finite family $(Aj)_{j\in J}$ of orbits of
    $\gal(\sqrt[n]{1}/F)=H\subset (\dZ/n)^\times$ in $\dZ/n$. We suppose that
    $A_j\ne\{0\}$ and that $\sum_{j\in J} \sum_{a\in A_j} a=0$.
  \item We take $I=\sqcup_{j\in J} A_j$; Galois acts on each $A_j$.
  \item We take $\bar F\supset \dQ(\sqrt[n]{1})$, $k=\dQ(\sqrt[n]{1})$ and, if
    $i\in I$ corresponds to $a\in A_j\subset (\dZ/n)^\times$,
    $_n\lambda_i(x)=x^a$.
\end{enumerate}

One checks using \ref{VI:4-20} that a general $j_\mu$ (relative to
$F_1$) is deduced from a $j_\mu$ as above (for
$F\subset \dQ(\sqrt[n]{1})$, $F_1$ an extension of $F$) by restriction to
$\gal(\bar F/F_1)$ and multiplication by a character of order two.

For $a\in \gal(k/\dQ)=(\dZ/n)^\times$, $a$ transforms the family of $A_j$
into that of the $a A_j$. If $a\in H$, the family of $A_j$ is therefore preserved,
and the $j_\mu(F_\fp)$ are invariant under $H$: 





\begin{proposition_}\label{VI:6-10}
Under the hypotheses of \ref{VI:6-9}, $j_\mu$ is defined by an algebraic
Hecke character of $F$ with values in $F$.
\end{proposition_}

Put again $F=\dQ(\sqrt[n]{1})$, and consider the Hecke
characters $j[a]$ defined by a family $(a_i)_{i\in I}$, $i\in \dZ/n$,
$a_i\ne 0$, $\sum a_i=0$. The following theorem was obtained
independently by Vishik. 





\begin{theorem_}\label{VI:6-11}
Every algebraic Hecke character of $\dQ(\sqrt[n]{1})$ is, up to a power,
in the group generated by the $j[a]$ and the norm.
\end{theorem_}

Since we work ``up to torsion,'' only the algebraic part
of the Hecke characters considered matters. Applying \eqref{VI:eq:6-8-1} to
the $j[a]$, with $a=$ an element of $\dZ/n$ repeated $n$ times, we
reduce to the following lemma.





\begin{lemma_}\label{VI:6-12}
Denote by $\sum x_j \sigma_j$ the elements of the group algebra
$\dQ[(\dZ/n)^\times]$, and let $X$ be the subgroup of those with
$x_j+x_{-j}=C^{\underline{t e}}$. Then $X$ is generated over $\dQ$ by
$N=\sum \sigma_j$ and by the
$g_a=\sum \left\lfloor \frac{i a}{n}\right\rfloor \sigma)i^{-1}$
($a\in \dZ/n$, $a\ne 0$).
\end{lemma_}

Both $X$ and $Y=\langle N,(g_a)\rangle$ are ideals of the group
algebra, and it suffices to check that they are annihilated by the same
complex characters $\chi$ of $(\dZ/n)^\times$. This amounts to saying that an
odd character $(\chi(j)=-\chi(-j)$) never annihilates all the $g_a$. If
$\chi$ is primitive, we have 
\[
  \chi(g_1) = \sum \frac i n \chi^{-1}(i) = -L(\chi,0)\ne 0 \text{.} 
\]
In the general case, if $\chi$ comes from a primitive character of
$(\dZ/n)^\times$, $\chi(g_{n/d})$ reduces to an analogous sum over
$(\dZ/n)^\times$, whence the lemma. 










\section{Generalised Kloosterman sums}\label{VI:7}





\subsection{}\label{VI:7-1}

Let $\dF_q$ be a finite field with $q$ elements, and $\psi$ a
non-trivial additive character. In this section, we make a cohomological study
of generalised Kloosterman sums 
\begin{equation*}\tag{7.1.1}\label{VI;eq:7-1-1}
  \kloos_{n,a} = \sum_{x_1\dotsm x_n = a} \psi(x_1+\cdots + x_n) \qquad (n\geqslant 1) \text{.} 
\end{equation*}
For $a=0$, this sum is elementary (cf. \ref{VI:7-7}): 
\begin{equation*}\tag{7.1.2}\label{VI:eq:7-1-2}
  \kloos_{n,0} = (-1)^{n-1} \text{.} 
\end{equation*}
The interesting case is that where $a\ne 0$. The $x_i$ are then in
$\dF_q^\times$. We will prove in this case a bound 
\begin{equation*}\tag{7.1.3}\label{VI:eq:7-1-3}
  \left|K_{n,a}\right| \leqslant n q^{\frac{n-1}{2}} \text{.} 
\end{equation*}

Our essential tools will be the cohomological statements reflecting
the following identities 
\begin{align*}\tag{7.1.4}\label{VI:eq:7-1-4}
  \kloos_{n,a} &= \sum_{x_1} \psi(x_i) \sum_{x_2\dotsm x_n = a/x_1} \psi(x_2+\cdots + x_n) \\ \notag
    &= \sum_{x\in \dF_q^\times} \psi(x) \kloos_{n-1,a/x} \qquad (a\ne 0,n\geqslant 2) \\ \tag{7.1.5}\label{VI:eq:7-1-5}
  \sum_a \kloos_{n,a} &= \sum_a \sum_{x_1 \dotsm x_n = a} \psi\left(\sum x_i\right) = \sum_{x_1,\dots,x_n} \psi\left(\sum x_i\right) = 0 
\end{align*}
For $\chi$ a non-trivial character of $\dF_q^\times$, 
\begin{equation*}\tag{7.1.6}\label{VI:eq:7-1-6}
  \sum \chi(a) \kloos_{n,a} = \sum_{x_1,\dots,x_n} \chi\left(\prod x_i\right) \cdot \psi\left(\sum x_i\right) = (-\tau(\chi^{-1},\psi))^n 
\end{equation*}
(Gauss sum). 





\subsection{}\label{VI:7-2}

Let $k$ be an étale algebra of degree $n$ over $\dF_q$, and $\varepsilon(k)$
as in \eqref{VI:eq:4-5-1}. Put 
\[
  \kloos_{k,a} = \sum_{N_{k/\dF_q}(x)=a} \psi \tr(x) \text{.} 
\]
We still have 
\begin{align*}\tag{7.2.1}\label{VI:eq:7-2-1}
  \kloos_{k,0} &= (-1)^{n-1} \cdot \varepsilon(k) \\ \tag{7.2.2}\label{VI:eq:7-2-2}
  \sum_a \kloos_{k,a} &= 0 \\ \tag{7.2.3}\label{VI:eq:7-2-3}
  \sum_a \chi(a) \kloos_{k,a} &= (-1)^n \tau_{\dF_q}\left(\chi^{-1}\circ N_{k/\dF_q},\psi\right) \text{.} 
\end{align*}

By Fourier inversion on $\dF_q^\times$, we deduce from this an expression for
$\kloos_{k,a}$ in terms of Gauss sums: 
\begin{align*} 
  \kloos_{k,a} &= \frac{1}{q-1} \sum_\chi \chi(a) \sum_{b\in \dF_q^\times} \chi^{-1}(b) \kloos_{k,b} \\ \tag{7.2.4}\label{VI:eq:7-2-4}
  (-1)^n \kloos_{k,a} &= \frac{1}{q-1} \left(\varepsilon(k) + \sum_{\chi\ne 1} \chi(a) \tau_{\dF_q}\left(\chi\circ N_{k/\dF_q},\psi\right) \right) \text{.} 
\end{align*}

The Hasse-Davenport identity now makes it possible to compare $\kloos_{k,a}$
with $\kloos_{n,a} = \kloos_{\dF_{q^n},a}$: 
\begin{equation*}\tag{7.2.5}\label{VI:eq:7-2-5}
  \kloos_{k,a} = \varepsilon(k) \kloos_{n,a} \text{.} 
\end{equation*}

I know of no cohomological interpretation of \eqref{VI:eq:7-2-4}, but
I will give one of \eqref{VI:eq:7-2-5}. Let us return to the sums $\kloos_{n,a}$. 





\subsection{}\label{VI:7-3}

Let $\dF$ be an algebraic closure of $\dF_q$ and, for $a\in \dF$, let
$V_a^{n-1}$, or simply $V_a$, be the hypersurface of $\dA^n$ of equation
$x_1 \dotsm x_n=a$. As usual, we will regard $\psi$ as taking
$\ell$-adic rather than complex values. With notation
\ref{VI:1-8}(ii), \eqref{VI:eq:7-1-3} follows from the 





\begin{theorem_}\label{VI:7-4}
The cohomology of $V_a^{n-1}$ with values in $\sF(\psi(\sum x_i))$ satisfies
\begin{enumerate}[\indent (i)]
  \item $\h_c^i=0$ for $i\ne n-1$;
  \item $\h_c^\bullet{} \iso \h^\bullet{}$;
  \item for $a\ne 0$, $\dim \h_c^{n-1}{}=n$;
  \item for $a=0$, $\h_c^{n-1}{}$ is canonically isomorphic to $E_\lambda$.
\end{enumerate}
\end{theorem_}





\subsection{}\label{VI:7-5}

As usual, this theorem also controls the dependence of
$\kloos_{n,a}$ on $q$: for each $a\in \dF_q^\times$, there exist $n$
Frobenius eigenvalues $\alpha_1,\dots,\alpha_n$, of complex absolute value
$q^{\frac{n-1}{2}}$, such that 
\[
  \sum_{\substack{x_1\dotsm x_n = a \\ x_i\in \dF_{q^m}}} \psi\circ \tr\left(\sum x_i\right) = (-1)^{n-1} \left(\alpha_1^m + \cdots + \alpha_n^m \right) \text{.} 
\]

For $n$ even, $x\mapsto -x$ is an involution of $V_a$, which transforms
$\sF(\psi)$ into its dual. Arguing as in \ref{VI:3-6}, we deduce from this a
canonical non-degenerate alternating form on $\h_c^{n-1}(V_a,\sF(\psi))$,
with values in $\dQ_\ell(-(n-1))$ and the 





\begin{corollary_}\label{VI:7-6}
With the notations of \ref{VI:7-5}, if $n$ is even, the $\alpha_i$ group
into $\frac n 2$ pairs of roots $\alpha,\bar\alpha$ of product equal to
$q^{n-1}$.
\end{corollary_}





\subsection{}\label{VI:7-7}

Let us check the case $a=0$ of \ref{VI:7-5}. The hypersurface $V_0$ is the union
of the coordinate hyperplanes $x_i=0$. Let us compute the Leray spectral
sequence \eqref{VI:eq:2-6-1*} of this covering of $V_0$, for cohomology with or
without support with coefficients in $\sF(\psi)$. By \ref{VI:2-7*}*, all
initial terms vanish, except the cohomology of the intersection $\{0\}$ of
all coordinate hyperplanes. There remains an isomorphism
$\h_c^0{}\iso \h^0{}=E_\lambda$ which justifies \ref{VI:7-4}.

We will prove \ref{VI:7-4} and \ref{VI:7-8} below by a simultaneous induction
on $n$, relying on the spectral sequence of which
\eqref{VI:eq:7-1-4} is the reflection. This requires controlling the dependence
on $a$ of $\kloos_{n,a}$.

Let $\pi:\dA^n \to \dA$ be the ``product of coordinates'' map, and
$\sigma$ the sum of coordinates. The hypersurfaces $V_a$ are the fibres of
$\pi$. 





\begin{theorem_}\label{VI:7-8}
\begin{enumerate}[(i)]
  \item The sheaf $\eR^{n-1}\pi_! \sF(\psi\sigma)$ is smooth of rank $n$ on
    $\dA^1\setminus \{0\}$.
  \item its extension by $0$ to $\dP^1$ is the direct image of its
    restriction to $\dA^1\setminus \{0\}$.
  \item at $0$, the monodromy is unipotent, with a single Jordan block
  \item at $\infty$, wild inertia acts without non-zero fixed point, and the
    Swan conductor is $1$,
  \item we have $\eR^i \pi_! \sF(\psi\sigma) \iso \eR^i \pi_\ast \sF(\psi\sigma)$
    (zero for $i\ne n-1$).
\end{enumerate}
\end{theorem_}
\begin{proof}[Proof]
That \ref{VI:7-4} (for a given value of $n$) implies \ref{VI:7-8}
(for the same value of $n$).

The fibre of $\eR^i \pi_! \sF(\psi\sigma)$ at $a$ is
$\h_c^i(V_a,\sF(\psi\sigma))$, and over the dense open of $\dA^1$ where the
$\eR^i \pi_!\sF(\psi^{-1}\sigma)$ are locally constant that of
$\eR^i \pi_\ast \sF(\psi\sigma)$ is $\h^i(V_a,\sF(\psi\sigma))$
(\hyperref[VII]{Th.\ Finitude} \ref{VII:2-1}). Hypothesis \ref{VI:7-4}($n$)
therefore gives us the 





\begin{lemma_}\label{VI:7-9}
\begin{enumerate}[(i)]
  \item $\eR^i \pi_! \sF(\psi\sigma)=0$ for $i\ne n-1$.
  \item For $i=n-1$, the fibre of this sheaf at every point $a\ne 0$ is of
    constant rank $n$. At $0$, it is of rank $1$.
  \item Over a dense open $U$, we have
    $\eR^i \pi_! \sF(\psi) \iso \eR^i \pi_\ast \sF(\psi)$.
\end{enumerate}
\end{lemma_}





\subsection{}\label{VI:7-10}

We will now use the Leray spectral sequence of $\pi$, for
cohomology with or without support, with coefficients in $\sF(\psi\sigma)$.
The abutment is known:
$\h_c^\bullet(\dA^n,\sF(\psi\sigma))=\h^\bullet(\dA^n,\sF(\psi\sigma)) = 0$
(\ref{VI:2-7*}*). In proper-support cohomology, \ref{VI:7-9}(i) ensures that
$E_2^{pq}=0$ for $q\ne n-1$. The $E_2$-terms therefore all vanish: 
\begin{equation*}\tag{7.10.1}\label{VI:eq:7-10-1}
  \h_c^p(\dA^1,\eR^{n-1} \pi_!\sF(\psi\sigma)) = 0 \text{.} 
\end{equation*}

For $p=0$, this means that $\eR^{n-1}\pi+! \sF(\psi\sigma)$ has no section
with punctual support. In view of the rank constancy \ref{VI:7-9}(ii), we deduce 
\begin{equation*}\tag{7.10.2}\label{VI:eq:7-10-2}
  \begin{array}{c}
    \text{The sheaf $\eR^{n-1} \pi_! \sF(\psi\sigma)$ on $\dA^1$ is smooth on $\dA^1\setminus \{0\}$, and is} \\
    \text{a subsheaf of the direct image $\sG$ of its restriction to $\dA^1\setminus \{0\}$.}
  \end{array}
\end{equation*}

The long exact cohomology sequence of 
\[\xymatrix{
  0 \ar[r] 
    & \eR^{n-1} \pi_! \sF(\psi\sigma) \ar[r] 
    & \sG \ar[r] 
    & \sQ \ar[r] 
    & 0 \text{.} 
}\]
where the support of $\sQ$ is concentrated at $0$, gives 
\[\xymatrix{
  \h_c^0(\dA^1,\sG) \ar[r] 
    & \sQ_0 \ar[r] 
    & \h_c^1(\dA^1,\eR^{n-1} \pi_! \sF(\psi\sigma)) \text{.} 
}\]
The extreme terms being zero, $\sQ=0$ and 
\begin{equation*}\tag{7.10.3}\label{VI:eq:7-10-3}
  \text{$\eR^{n-1} \pi_! \sF(\psi\sigma)$ is the direct image of its restriction to $\dA^1\setminus \{0\}$.}
\end{equation*}

Let $j$ be the inclusion of $\dA^1$ into $\dP^1$. The ``derived
category'' version of the Leray spectral sequences is 
\begin{align*}
  \hh^\bullet(\dP^1,j_! \eR\pi_! \sF(\psi\sigma)) &= \hh_c^\bullet(\dA^1,\eR \pi_! \sF(\psi\sigma)) = \h_c^\bullet(\dA^n,\sF(\psi\sigma)) = 0 \\
  \hh^\bullet(\dP^1,\eR j_\ast \eR \pi_\ast \sF(\psi\sigma)) &= \hh^\bullet(\dA^1,\eR \pi_\ast \sF(\psi\sigma)) = \h^\bullet(\dA^n,\sF(\psi\sigma)) = 0 \text{.}
\end{align*}
Let $\Delta$ be the mapping cylinder of
$j_! \eR\pi_! \sF(\psi\sigma) \to \eR j_\ast \eR \pi_\ast \sF(\psi\sigma)$.
By \ref{VI:7-9}(iii), the cohomology sheaves of this complex of
sheaves have finite support. On the other hand, the long exact
cohomology sequence of the triangle
$(j_!\eR \pi_! \sF(\psi\sigma),\eR j_\ast \eR \pi_\ast \sF(\psi\sigma),\Delta)$
shows that $\hh^\bullet(\dP^1,\Delta)=0$. We have
$\hh^\bullet(\dP^1,\Delta) = \h^0(\dP^1,\sH^\bullet(\Delta))$, and finally
$\Delta=0$: 
\[\xymatrix{
  j_! \eR\pi_! \sF(\psi\sigma) \ar[r]^-\sim 
    & \eR j_\ast \eR \pi_\ast \sF(\psi\sigma) \text{.} 
}\]

In particular, $\eR\pi_! \sF(\psi\sigma)\iso \eR\pi_\ast \sF(\psi\sigma)$ and,
these complexes of sheaves having only one non-zero cohomology sheaf, 
\begin{equation*}\tag{7.10.4}\label{VI:eq:7-10-4}
\xymatrix{
  j_! \eR \pi_! \sF(\psi\sigma) \ar[r]^-\sim 
    & j_\ast \eR \pi_! \sF(\psi\sigma) \text{.} 
}
\end{equation*}
This completes the proof of (i) (ii) (v). 





\subsection{}\label{VI:7-11}

Let $\swan_0$ and $\swan_\infty$ be the Swan conductors at $0$ and at
$\infty$ of $\eR^{n-1} \pi_! \sF(\psi\sigma)$. We have
$\chi_c(\eR^{n-1} \pi_! \sF(\psi\sigma))=0$ \eqref{VI:eq:7-10-1} and the
Euler-Poincaré formula \eqref{VI:eq:3-2-1} reduces to 
\begin{equation*}\tag{7.11.1}\label{VI:eq:7-11-1}
  \swan_0 + \swan_\infty =1 \text{:} 
\end{equation*}
one of these conductors is $0$ (tame ramification), the other is $1$. 





\subsection{}\label{VI:7-12}

Let $\sG$ be a sheaf of rank $1$ on $\dG_m$, of Kummer type
\ref{VI:4-7} ($\sG=\sK_n(\chi)$) and non-constant ($\chi\ne 1$). The sheaf
$\sF(\psi\sigma)\otimes \pi^\ast \sG$ on
$V^\ast=\pi^{-1}(\dG_m) = \dA^n\setminus \pi^{-1}(0)$ is then of the type
encountered in the study of Gauss sums, and \ref{VI:4-11}(ii) gives us 
\begin{equation*}\tag{7.12.1}\label{VI:eq:7-12-1}
\xymatrix{
  \h_c^\bullet(V^\ast,\pi^\ast\sG\otimes \sF(\psi\sigma)) \ar[r]^-\sim 
    & \h^\bullet(V^\ast,\pi^\ast \sG\otimes \sF(\psi\sigma)) \text{.} 
}
\end{equation*}
The Leray spectral sequence of $\pi$ transforms this statement into 
\begin{equation*}\tag{7.12.2}\label{VI:eq:7-12-2}
\xymatrix{
  \h_c^\bullet(\dG_m,\sG\otimes \eR^{n-1} \pi_! \sF(\psi\sigma)) \ar[r]^-\sim 
    & \h^\bullet(\dG_m,\sG\otimes \eR^{n-1} \pi_\ast \sF(\psi\sigma)) \text{.} 
}
\end{equation*}

Let $i$ be the inclusion of $\dG_m$ into $\dP^1$, and $\Delta$ the mapping cylinder
of
$\delta:i_!(\sG\otimes \eR^{n-1}\pi_! \sF(\psi))\to \eR i_\ast (\sG\otimes \eR^{n-1} \pi_! \sF(\psi\sigma))$
(where $\eR^{n-1} \pi_! = \eR^{n-1}\pi_\ast$). Arguing as in
\ref{VI:7-10}, we deduce from \eqref{VI:eq:7-12-2} that $\Delta=0$. In
particular 
\begin{equation*}\tag{7.12.3}\label{VI:eq:7-12-3}
  \text{At $0$ and at $\infty$, inertia acts without non-zero fixed point on $\sG\otimes \eR^{n-1}\pi_! \sF(\psi\sigma)$.}
\end{equation*}

At $\infty$, this holds even if $\sG$ is trivial \eqref{VI:eq:7-10-4}, so
that wild inertia acts without fixed point. In particular, the
ramification is wild: by \ref{VI:7-11}, we have $\swan_\infty=1$, and the
ramification at $0$ is tame.

At $0$, \eqref{VI:eq:7-12-3} forces the -- tame -- ramification of
$\eR^{n-1} \pi_! \sF(\psi\sigma)$ to be unipotent. By
\eqref{VI:eq:7-10-3} and \ref{VI:7-9}(ii), we have a single Jordan block. This
completes the proof of \ref{VI:7-8}($n$). 





\subsection{}\label{VI:7-13}

For $n=1$, \ref{VI:7-4} is trivial. To complete the proof of \ref{VI:7-4} and
\ref{VI:7-8}, it remains for us to prove that \ref{VI:7-8}, for a given
value of $n$, implies \ref{VI:7-4}, for $n+1$. The case $a=0$ having been
treated (\ref{VI:7-7}), we suppose that $a\ne 0$. We will write $x_0,\dots,x_n$ for
the coordinates of $\dA^{n+1}$. Let $g=x_0:V_a\to \dG_m$, $\tau$
the involution $x\mapsto a x^{-1}$ of $\dG_m$, and let as above $\pi$ be the
product of coordinates: $\dA^{n+1} \to \dA^1$, and $\sigma$ their sum. We
will denote by $\sF(\psi\sigma)$ the sheaf $\sF(\psi(\sum x_i))$ both on $\dA^{n+1}$
and on $\dA^n$. 





\subsection{}\label{VI:7-14}

Let us write the Leray spectral sequence of $g$. The sheaf $\sF(\psi\sigma)$ on
$\dA^{n+1}$ being the external tensor product of the sheaves $\sF(\psi)$ on
$\dA^1$ and $\sF(\psi\sigma)$ on $\dA^n$, we find (cf. \eqref{VI:eq:7-1-4}). 
\begin{align*}
  'E_2^{p q} &= \h_c^p(\dG_m,\sF(\psi)\otimes \tau^\ast \eR^q \pi_! \sF(\psi\sigma)) \Rightarrow \h_c^{p+q}(V_a,\sF(\psi\sigma)) \\
  ''E_2^{p q} &= \h^p(\dG_m,\sF(\psi)\otimes \tau^\ast \eR^q \pi_\ast \sF(\psi\sigma)) \Rightarrow \h^{p+q}(V_a,\sF(\psi\sigma)) \text{.} 
\end{align*}
By hypothesis,
$\eR^q \pi_! \sF(\psi\sigma) \iso \eR^q \pi_\ast \sF(\psi\sigma)$, and this
sheaf vanishes for $q\ne n-1$. The sheaf $\sF(\psi)$ on $\dG_m$ is
wildly ramified at $\infty$ and unramified at $0$, while
$\tau^\ast \eR^q \pi_! \sF(\psi\sigma)$ is wildly ramified at $0$ (without
invariant under wild inertia, of Swan conductor $1$) and tamely
ramified at $\infty$. Their tensor product is therefore
\begin{enumerate}[\indent a)]
  \item wildly ramified at $0$ and $\infty$, without invariant under wild
    inertia;
  \item of Swan conductor $1$ at $0$ and $n$ (the rank of
    $\eR^q \pi_! \sF(\psi))$ at $\infty$.
\end{enumerate}
We have $'E_2^{p q} \iso ''E_2^{p q}$, and $'E_2^{p q}=0$ except for $p=1$ and
$q=n-1$. The Euler-Poincaré formula finally gives 
\[
  \dim{'E_2^{p q}} = 1+n \text{;} 
\]
this completes the proof.
\end{proof}





\subsection{Remark}\label{VI:7-15}

The isomorphism obtained in \ref{VI:7-14} 
\begin{equation*}\tag{7.15.1}\label{VI:eq:7-15-1}
  \h_c^n(V_a^n,\sF(\psi\sigma)) = \h_c^1(\dG_m,\sF(\psi)\otimes \tau^\ast \eR^{n-1} \pi_! \sF(\psi\sigma)) 
\end{equation*}
allows, with the notations of \ref{VI:7-5}, computing the product $\det(F)$
of the Frobenius eigenvalues $\alpha_i$. The formalism of \cite{de73}
applies, since $\sF(\psi)$ and $\eR^{n-1} \pi_! \sF(\psi\sigma)$
belong to infinite compatible systems of $\ell$-adic
representations. This reduces us to local problems, which one can
simplify by noting that (cf. \cite[9.5]{de73})
\begin{enumerate}[\indent a)]
  \item the global constants for
    $\sF(\psi)$ and $\tau^\ast \eR^{n-1} \pi_! \sF(\psi)$ are easy to
    compute, the cohomology of these sheaves being of trivial nature.
  \item at every point, for one or the other of these sheaves, the
    corresponding $\ell$-adic representation of the decomposition
    group has unramified semi-simplification.
\end{enumerate}

The result, for the sum in $n$ variables, is that 
\begin{equation*}\tag{7.15.2}\label{VI:eq:7-15-2}
  \det\left(F,\h_c^{n-1}(V_a^{n-1},\sF(\psi))\right) = q^{\frac{n(n-1)}{2}} \text{.} 
\end{equation*}





\subsection{Remark}\label{VI:7-16}

For $p=2$, the sheaf $\sf(\psi)$ is orthogonal. For $a\ne 0$ and $n$
odd, we deduce by \ref{VI:7-4}(ii) and Poincaré duality a
symmetric bilinear form on $\h_c^{n-1}(V_a^{n-1},\sF(\psi))$, with
values in $\bar\dQ_\ell(1-n)$. Relative to this form, $F$ is an
orthogonal similitude of multiplier $q^{n-1}$, and $q^{\frac{1-n}{2}} F$
lies in the special orthogonal group, by \eqref{VI:eq:7-15-2}. One
of the eigenvalues of $F$ is therefore $q^{\frac{1-n}{2}}$, and the others fall
into pairs of roots $\alpha$, $\bar\alpha$ of product $q^{n-1}$.

For $n=3$, L.\ Carlitz \cite{ca69-2} obtained a more precise result,
equivalent to the following proposition. 





\begin{proposition_}\label{VI:7-17}
For $p=2$, and $\psi(x)=(-1)^{\tr_{\dF_q/\dF_q}(x)}$, we have an isomorphism 
\[
  \h_c^2(V_a^2,\sF(\psi\sigma)) = \operatorname{Sym}^2 \h_c^1(V_a^1,\sF(\psi\sigma)) \text{.} 
\]
\end{proposition_}

The second tensor power of $\h_c^1(V_a^1,\sF(\psi\sigma))$ is
$\h_c^1(V_a^1\times V_a^1,\sF(\psi\sigma)\boxtimes \sF(\psi\sigma))$
(K\"unneth), the symmetry $x\otimes y\mapsto y\otimes x$ being represented
geometrically as $-\tau^\ast$, where $\tau$ is the automorphism
$(x,y)\mapsto (y,x)$ of $V_a\times V_a$. If $X'$ is the quotient of
$V_a^1\times V_a^1$ by $\{\operatorname{id}{},\tau\}$ and $\pi$ is the
projection from $V_a^1\times V_a^1$ to $X'$, the group
$\operatorname{Sym}^2 \h_c^1(V_a^1,\sF(\psi\sigma))$ therefore identifies with the
cohomology of $X'$ with coefficients in the $\tau$-anti-invariant part of
$\pi_\ast(\sF(\psi\sigma)\boxtimes \sF(\psi\sigma))$. Denote by $s$ the function
on $X'$ with $s\pi =\sigma\pr_1+\sigma\pr_2$. We have 
$\sF(\psi\sigma)\boxtimes \sF(\psi\sigma) = \sF(\psi(\sigma\pr_1+\sigma\pr_2)) = \sF(\psi s \pi) = \pi^\ast \sF(\psi s)$ 
(isomorphism compatible with $\tau$). On the image $\pi(\Delta)$ of the
diagonal, the sought anti-invariant sheaf is therefore zero, while on
the complement $X=X'\setminus \pi(\Delta)$, it is the tensor product of
$\sF(\psi s)$ with $\varepsilon(P)$, for $\varepsilon$ the unique character
of order $2$ of $\{1,\tau\}$, and $P$ the $\{1,\tau\}$-torsor over $X$ that is the
double covering induced by $V_a^1\times V_a^1$. We have 
\begin{equation*}\tag{7.17.1}\label{VI:eq:7-17-1}
  \operatorname{Sym}^2 \h_c^1(V_a^1,\sF(\psi\sigma)) = \h_c^2(X,\sF(\psi\sigma)\otimes \varepsilon(P)) \text{.}
\end{equation*}

Let us compute $X$, $s$ and $\varepsilon(P)$. On $V_a^1\times V_a^1$, put
$x_1=x_1\circ\pr_1$, $x_2=x_2\circ\pr_2$, $x_1'=x_1\circ \pr_2$,
$x_2'=x_2\circ \pr_2$. We have $x_1 x_2=x_1'x_2'=a$. Let us identify functions on
$X'$ with $\tau$-invariant functions on $X$. Let $s_1=x_1+x_1'$,
$s_2=x_2+x_2'$, $p_1=x_1 x_1'$, $p_2=x_2 x_2'$. We have $s_1 p_2 = a s_2$,
$s_2 p_1 = a s_1$ and $X$ is defined in $X'$ by $s_1\ne 0$, $s_2\ne 0$. On
$X$, we therefore have $p_1 = a s_1 s_2^{-1}$, and $(s_1,s_2)$ identifies $X$ with
$\dG_m\times \dG_m$. On $X$, the covering $P$ has equation
$T^2-s_1 T+a s_1 s_2^{-1}$ (take for $T$ the coordinate $x_1$). If
$T=s_1 t$, this is rewritten $t^2-t+a s_1^{-1} s_2^{-1}$. We therefore have
$\varepsilon(P) = \sF(\psi(a s_1^{-1} s_2^{-1}))$. Since $s=s_1+s-2$, we have
$\sF(\psi s)\otimes \varepsilon(P) = \sF(\psi(s_1+s_2+ a s_1^{-1} s_2^{-1}))$
and the right-hand side of \eqref{VI:eq:7-17-1} identifies with the left-hand side of
\ref{VI:7-17}. 


\paragraph{Remark}

$\h_c^1(V_a^1,\sF(\psi\sigma))$ is also the first cohomology group of
the completed elliptic curve of the affine curve of equation 
\begin{align*}
  y^2-y &= x_1+x_2 \\
  x_1 x_2 &= a \text{;} 
\end{align*}
its modular invariant is $j=a^{-2}$. 





\subsection{Remark}\label{VI:7-18}

For $\chi$ a multiplicative character of $\dF_q^\times$, and $a\ne 0$, the
method of \ref{VI:7-14} applies to the study of the sum
\[
  s = \sum_{x_0\dotsm x_n = a} \chi(x_0) \psi(x_0+\cdots + x_n)
\]
(replace the sheaf $\sF(\psi)$ on $\dG_m$ by $\sF(\chi\cdot \psi)$). We
still find a cohomology of dimension $n+1$, and
$|S|\leqslant (n+1)q^{n/2}$. Such sums have been considered by
Salie and Mordell \cite{mo73}.

One may hope that our methods make it possible to study the sums
\[
  S_{k,\chi,a} = \sum_{N(x)=a} \chi(x) \psi \tr(x)
\]
($k$ an étale algebra over $\dF_q$, $a\in \dF_q^\times$, $\chi$ a character of
$k^\times$). The analogue of \ref{VI:7-4} (for $a\ne 0$) should hold. 





\subsection{}\label{VI:7-19}

The symmetric group $S_n$ acts on $\dA^n$ preserving $V_a$ and
the map $\sigma$, hence also the sheaf $\sf(\psi\sigma)$, inverse
image by $\sigma$ of the sheaf $\sf(\psi)$ on $\dG_a$. It therefore acts
on $\h_c^{n-1}(V_a,\sF(\psi\sigma))$. Formula \eqref{VI:eq:7-2-5} admits
the following cohomological interpretation





\begin{proposition_}\label{VI:7-20}
The action of $S_n$ on $\h_c^{n-1}(V_a,\sF(\psi\sigma))$ is multiplication
by the sign character.
\end{proposition_}

This amounts to showing that a transposition $\tau$ acts by multiplication
by $-1$. Since $\tau^2=1$, $\tau$ can only have $\pm 1$ as eigenvalues.
Let $V_a/\tau$ be the quotient of $V_a$ by
$\{\operatorname{id},\tau\}$. The map $\sigma$, hence the sheaf
$\sf(\psi\sigma)$, pass to the quotient, and the subspace of
$\h_c^{n-1}(V_a,\sF(\psi))$ fixed by
$\tau$ is $\h_c^{n-1}(V_a/\tau,\sF(\psi\sigma))$. We must prove that
this cohomology vanishes.

Suppose that $a\ne 0$. For $\tau=(1,2)$ the quotient $V_a/\tau$ identifies
with $\dG_a\times \dG_m^{n-2}$, with quotient map
$(x_1,x_2,x_3,\dots,x_n)\mapsto (x_1+x_2,x_3,\dots,x_n)$. Indeed, $x_1$ and
$x_2$ are determined up to order by $x_1+x_2$ and
$x_1x_2 = a/x_3\dotsm x_n$. In these coordinates $(t,x_3,\dots,x_n)$, the
sheaf is written
$\sF(\psi(t+x_3+\cdots + x_n)) = \sF(\psi(t))\otimes \sF(\psi(x_3+\cdots + x_n))$.
The K\"unneth formula, and $\h_0^\bullet(\dG_a,\sF(\psi)) = 0$, therefore
give the required vanishing.

For $a=0$, one can deduce \ref{VI:7-20} from computation \ref{VI:7-7}. 





\subsection{}\label{VI:7-21}

The sums $\kloos_{n,a}$ of this section have been studied by
S.\ Sperber, in his thesis, and in \cite{sp77} by $p$-adic
methods inspired by Dwork and by Bombieri's article \cite{bo66}.
For $p\ne 2$, these methods give him the dependence on $q$ \ref{VI:7-5},
the functional equation $\alpha_i = q^{n-1} \beta_i^{-1}$ between the
Frobenius eigenvalues for the sums $\kloos_{n,a}$ and
$\kloos_{n,b}$ if $b=(-1)^n a$ (cf. \ref{VI:7-6}), and the value \ref{VI:7-15}
of the product of the roots. For $p\geqslant n+3$, he moreover determines the
$p$-adic valuations of the roots. The result is very beautiful: if one orders
the roots $\alpha_i$ suitably, with
$0\leqslant i\leqslant n$, the $\alpha_i q^{-i}$ are units. The
restrictions on $p$ are no doubt not essential. 










\section{Other applications}\label{VI:8}





\subsection{}\label{VI:8-1}

Identity \ref{VI:1-9-4} can sometimes be used to compute
the number of rational points of an algebraic variety over $\dF_q$. In
most cases where an exact result has been obtained, cohomology
is expressed in terms of algebraic cycles; such is the case for rational
surfaces, quadrics, smooth intersections of two quadrics of
odd dimension, flag varieties\ldots. This is however not
always the case; in Lusztig's very beautiful article \cite{lu76}, algebraic
cycles do not appear. 





\subsection{}\label{VI:8-2}

For reductive groups, cohomology can be computed by lifting to
characteristic $0$: if $B$ is a Borel subgroup of $G$, $G$ is a
$B$-torsor over $G/B$; the cohomology of the smooth projective variety
$G/B$ is invariant under specialisation, and we use the Leray
spectral sequence of $G\to G/B$ to prove the same result for $G$. It is
conveniently expressed in terms of \emph{the} maximal torus $T$ of $G$ and of the
action of \emph{the} Weil group $W$ on $T$ (for the definition of \emph{the}
maximal torus, cf. \cite[VIII \S 2 Rem.2]{bo05} or \cite[p.105]{dl76}). It is
the exterior algebra of its primitive part, which coincides up to a shift
with the quotient $I(\h^\bullet(B G,\dQ_\ell))$ of the cohomology of
$B G$ formed by ``indecomposable elements,'' and
$\h^\bullet(B G,\dQ_\ell)=\h^\bullet(B T,\dQ_\ell)^W$. If $X$ is the
character group of $T$, we have in total 
\[
  \h^\bullet(G,\dQ_\ell) = \bigwedge \left(I\left(\operatorname{Sym}^\bullet(X\otimes \dQ_\ell(-1))^W\right)[-1]\right) \text{.}
\]
If $G$ is defined over a finite field $\dF_q$, $\gal(\dF/\dF_q)$ acts on $X$
and this isomorphism is Galois-compatible. Passing from there to
proper-support cohomology, we find the classical formulas for the
number of rational points of $G$. If $F^\ast:X\to X$ is defined by the
Frobenius morphism $F:T\to T$, and $\deg_G(F)$ denotes the degree
$q^{\dim G}$ of $F:G\to G$, we have 
\begin{equation*}\tag{8.2.1}\label{VI:eq:8-2-1}
  \left|G_0(\dF_q)\right| = \deg_G(F) \cdot \det\left(1-{F^\ast}^{-1},I(\operatorname{Sym}^\bullet(X\otimes \dQ)^W)\right) \text{.}
\end{equation*}

The same formula holds for Ree and Suzuki groups; these groups are
of the form $G^F$ ($G$ reductive, $F^2$ a Frobenius), and it suffices to
check that the Lefschetz trace formula holds for $F$ acting
on $G$. If $B$ is an $F$-stable Borel subgroup, one checks it
directly for $F$ acting on $B$, and for the action on $G/B$, proper and
smooth, one may invoke the general theorems. 





\subsection{}\label{VI:8-3}

The trace formula has been used by Deligne-Lusztig \cite{dl76},
Kazhdan \cite{ka77} and Springer \cite{sp76} in the study of
complex representations of the finite groups $G_0(\dF_q)$, for $G_0$
reductive over $\dF_q$. The works of Kazhdan and Springer contain
admirable examples of how the trace formula allows the ``analytic
continuation'' from a split situation to a non-split situation
(cf. \ref{VI:1-13}). 






